% !TeX program = xelatex
% Standalone continuation to be placed after Chapter 11 of the textbook.
% The file also compiles by itself in Overleaf with XeLaTeX.
\documentclass[11pt,letterpaper,openany]{book}

\usepackage[margin=1in]{geometry}
\usepackage{fontspec}
\setmainfont{Latin Modern Roman}
\setsansfont{Latin Modern Sans}
\setmonofont{Latin Modern Mono}
\usepackage{amsmath,amssymb}
\usepackage{booktabs,longtable,tabularx,array}
\usepackage{fancyvrb}
\usepackage{hyperref}
\usepackage{xcolor}

\hypersetup{
  colorlinks=false,
  pdfborder={0 0 0},
  pdftitle={Interview Preparation and Quantitative Finance Interview Questions}
}

\setlength{\parindent}{15pt}
\setlength{\parskip}{0pt}
\raggedbottom

\makeatletter
\def\ps@textbookdraft{%
  \def\@oddhead{\small DRAFT\hfil}%
  \def\@evenhead{\small DRAFT\hfil}%
  \def\@oddfoot{\hfil\thepage\hfil}%
  \def\@evenfoot{\hfil\thepage\hfil}%
}
\makeatother
\pagestyle{textbookdraft}

\DefineVerbatimEnvironment{codeblock}{Verbatim}{
  fontsize=\small,
  frame=single,
  framesep=2mm
}

\newcolumntype{L}[1]{>{\raggedright\arraybackslash}p{#1}}

\newcommand{\ChapterIntro}[1]{%
  \begin{quote}\small #1\end{quote}%
}

\newenvironment{EditorialNote}
  {\begin{quote}\small\setlength{\parskip}{0.25em}}
  {\end{quote}}

\newcommand{\QuestionSpacer}{\vspace{0.4em}}

\newcommand{\InterviewQuestionHeader}[5]{%
  \vspace{1.0em}%
  {\large\bfseries #1\enspace #4\par}%
  {\small\itshape Institution type: #2\par}%
  {\small\itshape Role: #3\enspace |\enspace Difficulty: #5\par}%
  \vspace{0.25em}\hrule\vspace{0.45em}%
}

\newcommand{\FoundationQuestionHeader}[3]{%
  \vspace{1.0em}%
  {\large\bfseries #1\enspace #2\par}%
  {\small\itshape Institution type: Cross-industry quantitative-finance employers\par}%
  {\small\itshape Category: #3\par}%
  \vspace{0.25em}\hrule\vspace{0.45em}%
}

\newcommand{\RoleProfile}[2]{%
  \subsection{#1}%
  \noindent\textit{Typical institution type: #2.}\par
}

\newcommand{\prompt}[1]{\subsubsection*{#1}}
\newcommand{\playbooknote}[1]{\begin{quote}\small #1\end{quote}}

\begin{document}
\setcounter{chapter}{11}
\chapter{Interview Preparation}

\section{Introduction}
Quantitative-finance interviews test more than isolated technical facts. A strong candidate must identify the type of institution and role, recognize the assessment format, structure an answer under time pressure, and connect analysis to a trading, investment, or risk decision. This chapter organizes preparation around those repeatable tasks.

Specific employer names and identifying details are intentionally omitted. Institution types are retained because they are useful predictors of interview emphasis. A proprietary trading or market-making firm will often emphasize probability, mental arithmetic, games, and market intuition; an investment bank may combine product knowledge, risk, coding, and behavioral assessment; an asset manager or institutional investor may place more weight on portfolio construction, research process, and communication; and a hedge fund or quantitative investment manager may emphasize modeling depth, statistics, and independent research judgment.

The most effective way to use this chapter is in three passes. First, learn the process maps and answer structures. Second, practice the cases and technical prompts without notes. Third, personalize every behavioral answer with evidence from projects, coursework, internships, or research.
\section{The End-to-End Interview Process}
The exact sequence varies, but most recruiting processes can be understood as a progression from broad screening to increasingly job-specific evidence. At each stage, the candidate should know three things: the format, the competency being tested, and the evidence needed to pass.

\begin{longtable}{@{}L{0.16\textwidth}L{0.25\textwidth}L{0.25\textwidth}L{0.26\textwidth}@{}}
\toprule
\textbf{Stage} & \textbf{What usually happens} & \textbf{What is being tested} & \textbf{How to prepare} \\
\midrule
\endfirsthead
\toprule
\textbf{Stage} & \textbf{What usually happens} & \textbf{What is being tested} & \textbf{How to prepare} \\
\midrule
\endhead
Role and application screen & R\'esum\'e review, eligibility checks, and sometimes a short application question. & Minimum technical fit, evidence of interest, and whether experience is relevant to the role. & Translate every major r\'esum\'e bullet into a problem, method, result, and lesson. Prepare a role-specific explanation of why the work matters. \\
Recruiter or HR screen & A short call covering motivation, logistics, communication, and broad experience. & Clarity, genuine interest, basic role understanding, and absence of obvious fit concerns. & Prepare 60-second answers for ``tell me about yourself,'' ``why this role,'' and ``why this institution type.'' Know dates, location constraints, and work authorization. \\
Online assessment & Timed mathematics, probability, coding, data analysis, pattern recognition, or numerical tests. & Speed, accuracy, fundamentals, and performance under constraints. & Rehearse in the actual time limit. Use a skip rule, track error types, and practice the platform mechanics before the scored attempt. \\
Pre-recorded video & Several behavioral or motivation questions with limited preparation and response time. & Structure, concision, presence, and consistency of motivation. & Practice 60--90 second answers on camera. Lead with the answer, use one example, and finish with relevance to the role. \\
First live interview & One or two interviews mixing behavioral questions with technical fundamentals. & Whether the candidate can reason interactively, accept hints, and explain work clearly. & Practice aloud with a partner. State assumptions, narrate the method, and check the result instead of silently calculating. \\
Take-home or case study & A coding assignment, research task, forecasting problem, investment case, or risk analysis, often followed by a presentation. & Problem definition, judgment, implementation quality, validation, and communication. & Define the target and evaluation metric first. Build a transparent baseline, document assumptions, test failure modes, and prepare a short executive summary. \\
Final round or superday & Several interviews with managers, technical specialists, and potential teammates. & Depth, consistency across interviewers, judgment, collaboration, and team fit. & Prepare for repeated questions without sounding scripted. Review every earlier submission and maintain the same assumptions, results, and motivation story. \\
Informal interaction & Lunch, coffee, office tour, or networking conversation. & Professionalism, curiosity, listening, and how the candidate behaves when the format feels less formal. & Ask thoughtful questions about work, feedback, and team decision-making. Treat the interaction as evaluated while remaining conversational. \\
Follow-up & Thank-you note, additional materials, references, or a clarification call. & Professional judgment and continued interest. & Send a concise note, correct any material error transparently, and avoid introducing new unsupported claims. \\
\bottomrule
\end{longtable}

\subsection{Typical Sequences by Institution Type}
Institution type is often a better predictor of the process than the exact job title.

\begin{itemize}
  \item \textbf{Proprietary trading or market-making firm:} timed mathematics or game assessment $\rightarrow$ probability and mental-arithmetic interview $\rightarrow$ trading games, market intuition, coding, and behavioral final rounds.
  \item \textbf{Hedge fund or quantitative investment manager:} mathematics or coding screen $\rightarrow$ statistics, machine learning, or research interview $\rightarrow$ take-home research task or deep project discussion $\rightarrow$ final interviews focused on judgment and research independence.
  \item \textbf{Investment bank or diversified financial institution:} recruiter or recorded-video screen $\rightarrow$ role-specific technical and behavioral interviews $\rightarrow$ multi-interviewer final round. Product knowledge, risk, controls, and communication may all appear.
  \item \textbf{Asset manager or institutional investor:} motivation and market discussion $\rightarrow$ investment, portfolio, or data case $\rightarrow$ presentation and panel interviews. The process often emphasizes research judgment and communication to stakeholders.
  \item \textbf{Insurance, risk, or advisory institution:} behavioral screen $\rightarrow$ statistics, financial-risk, product, or modeling interview $\rightarrow$ case study or manager panel. Governance and communication are often tested alongside technical knowledge.
\end{itemize}

\subsection{How to Read Signals During the Process}
The format itself tells the candidate how to allocate preparation time. A high-volume timed test rewards automaticity and a skip strategy. A take-home case rewards explicit assumptions, reproducibility, and documentation. A live technical interview rewards interaction: the interviewer needs to hear the reasoning and see how the candidate responds to hints. A final-round panel tests consistency; answers about motivation, projects, and technical choices should agree across interviewers.

\section{A Repeatable Preparation Method}

\subsection{Step 1: Decode the Role}
Reduce the job description to four columns: responsibilities, required knowledge, likely interview format, and personal evidence. Repeated verbs are especially informative. ``Research'' suggests hypothesis formation and validation; ``build'' suggests coding and systems design; ``manage risk'' suggests scenarios, limits, and communication; ``trade'' suggests probability, market intuition, and decisions under uncertainty.

\subsection{Step 2: Build a Competency Map}
For each role, score preparation needs from 0 to 3 across the following areas:
\begin{itemize}
  \item probability, statistics, and mathematical reasoning;
  \item coding, algorithms, SQL, and data manipulation;
  \item machine learning, time series, and validation;
  \item derivatives, fixed income, portfolio theory, or product knowledge;
  \item market intuition and commercial judgment;
  \item behavioral evidence, communication, and motivation.
\end{itemize}
Spend most practice time on high-importance areas where current performance is weak. Do not distribute time evenly simply because the syllabus is broad.

\subsection{Step 3: Build an Evidence Bank}
Prepare six to eight reusable stories rather than a separate story for every behavioral question. The bank should cover leadership, conflict, failure, ambiguity, rapid learning, ethical judgment, analytical impact, and communication with a non-technical stakeholder. Each story should contain a measurable result and a sentence connecting it to the target role.

\subsection{Step 4: Prepare Technical Answers at Two Depths}
Every core topic should have a 30-second explanation and a 90-second explanation. The short version gives the definition, intuition, and one limitation. The longer version adds a formula or example, implementation detail, and follow-up. This prevents both one-word answers and uncontrolled lectures.

\subsection{Step 5: Practice in the Real Format}
Untimed study builds knowledge; format-matched rehearsal builds interview performance. Complete timed assessments under the real limit, record video answers with the real preparation window, code without autocomplete when appropriate, present cases from slides, and solve live technical questions while speaking.

\subsection{Step 6: Keep an Error Log}
After every practice session, classify misses as concept, setup, arithmetic, coding, time management, or communication errors. Record the correction and one new drill. Improvement should be measured by repeated performance, not by whether the solution looked familiar after reading it.

\subsection{Preparation Deliverables}
Before the first interview, the candidate should have the following compact set of materials:
\begin{itemize}
  \item a one-page role and institution-type brief;
  \item a 60-second introduction and a 90-second motivation answer;
  \item six to eight behavioral stories;
  \item a project deep-dive sheet covering data, method, validation, result, limitations, and personal contribution;
  \item a technical formula and concept sheet tailored to the role;
  \item one timed-assessment strategy, including a skip rule;
  \item one mock case or coding exercise completed under realistic conditions;
  \item five thoughtful questions for interviewers at different seniority levels.
\end{itemize}
\section{Role-Specific Process Examples}

\RoleProfile{Commodities Trading (Summer Trading Analyst)}{Commodity merchant, energy trading firm, or commodities desk}
\textbf{Observed interview format}
\begin{itemize}
  \item Case study focused on natural gas price prediction (features $\to$ next-year price).
  \item HR interview (trading interest, commodity preference, data experience).
  \item Superday: in-person, multiple rounds, including manager behavioral interview.
  \item Networking lunch (informal but evaluated).
  \item Brain teasers / estimation questions.
  \item Market intuition: event-driven scenario (e.g., refinery explosion).
  \item Technical: Python + pandas, stochastic process simulation, Monte Carlo, plus SQL (medium-difficulty practice problems).
\end{itemize}

\RoleProfile{Quantitative Risk Analyst}{Cross-industry quantitative-finance employers}
\textbf{Core topics observed}
\begin{itemize}
  \item Value at Risk (VaR) and practical risk measurement.
  \item Data-driven risk modeling and stress testing.
  \item Delta--gamma option risk/trading questions.
\end{itemize}

\RoleProfile{Pre-Recorded Video Interview (pre-recorded video interview): Technology Audit Summer Analyst (Internal Audit)}{Investment bank or diversified financial institution}
\textbf{Format observed}
\begin{itemize}
  \item Pre-recorded video interview (pre-recorded video interview).
  \item Questions were primarily behavioral and experience-based.
  \item Emphasis on communication clarity and structured examples.
\end{itemize}

\textbf{Questions encountered}
\begin{enumerate}
  \item What interests you most about a career in Internal Audit?
  \item How do you see the Summer Analyst Program contributing to your long-term career goals?
  \item To what extent have you leveraged GenAI or other data analytics tools (Python, SQL, visualization and workflow-automation tools), and what impact will these technologies have on internal audit or financial services?
\end{enumerate}

\RoleProfile{Pre-Recorded Video Interview (pre-recorded video interview): Investment Management Division (Portfolio Solutions Group)}{Asset management firm or institutional investment division}
\textbf{Questions encountered}
\begin{enumerate}
  \item Why do you want to pursue a career in the Investment Management Division?
  \item Which group(s) within Investment Management are you most interested in and why?
  \item What is your greatest weakness in understanding finance/economics? What did you do to remediate it?
  \item ``Doing the right thing'' is a core value. Describe a time you noticed something that did not appear right. How did you handle it?
\end{enumerate}

\RoleProfile{Other Assessments Mentioned}{Multiple quantitative-finance institution types}
\textbf{Macro hedge fund (Risk and Treasury Intern)}
\begin{itemize}
  \item Cognitive screening assessment (12 minutes).
  \item Python assessment (about 60 minutes).
  \item Group exercise: 3--4 people working on a stock model.
\end{itemize}

\textbf{Other firms / platforms referenced}
\begin{itemize}
  \item Options market-making firm: online math assessment.
  \item Quant research: a timed mathematics assessment + cognitive section + brain teaser.
  \item Quantitative proprietary trading firm: online assessment.
  \item Electronic market-making firm: ``Beat the Odds'', number logic, and game-style tests (e.g., zap-N style tasks).
\end{itemize}

\RoleProfile{Quant Research Intern: Assessment Formats}{Hedge fund, quantitative investment manager, or proprietary trading firm}
\textbf{Track A: Proctored Math Online Assessment}
\begin{itemize}
  \item Proctored online math assessment.
  \item Duration: 90 minutes.
  \item Volume: 16 questions.
  \item Topics: probability, expected value, combinatorics, brain teasers.
\end{itemize}

\textbf{Track B: 24-Hour Take-Home Coding Assessment}
\begin{itemize}
  \item Take-home coding assessment with a 24-hour deadline.
  \item Content: object-oriented programming, understanding bid/ask spreads, arbitrage logic.
  \item Difficulty: conceptually easy, but long; completeness and code quality matter.
\end{itemize}

\textbf{Track C: Proctored Python Coding (Mostly pandas)}
\begin{itemize}
  \item Proctored Python coding assessment focused on pandas.
  \item Emphasis: EDA, data wrangling, modeling, evaluation.
  \item Difficulty: generally easy, but sometimes requires highly specific pandas details (e.g., \texttt{int64} dtype).
\end{itemize}

\RoleProfile{Trading Track Online Assessments}{Proprietary trading or electronic market-making firm}
\textbf{Track A: Long-Form Game-Theory Puzzle (Card Game Inference)}
\begin{itemize}
  \item A long puzzle session (about several hours).
  \item A card game was described with specific rules.
  \item Based on observed actions, the task was to infer which card(s) the person likely held.
  \item Required probability, expected value, and ``optimal play'' reasoning.
\end{itemize}

\textbf{Track B: Cognitive Ability Test (Pattern + Memory Games)}
\begin{itemize}
  \item Puzzle / memory game format.
  \item Pattern recognition: choose the next geometric sequence.
  \item Working memory: click the order of dots that blinked (sequence recall).
\end{itemize}

\textbf{Track C: Valuation and Expected Value Problems (EV/EBITDA + Default Risk)}
\begin{itemize}
  \item Finance/accounting flavored questions involving EV/EBITDA.
  \item Expected value problems for pricing assets under default probability or price decline scenarios.
  \item More straightforward calculations, but required comfort with accounting terms and valuation logic.
\end{itemize}

\RoleProfile{Trading Risk Management Intern}{Investment bank, market-making firm, hedge fund, or asset manager}
\textbf{Observed interview format}
\begin{itemize}
  \item Total of 3 interviews.
  \item Round 1: basic HR phone screen (qualification check; no technical).
  \item Round 2: 30 minutes with 2 managers (mixed technical + behavioral).
  \item Round 3: 30 minutes with 2 managers (mixed technical + behavioral).
\end{itemize}

\RoleProfile{Quantitative Technology Intern and Related Roles}{Cross-industry quantitative-finance employers}
\textbf{Track A: Quantitative Technology Intern (Three Interviewers)}
\begin{itemize}
  \item Round 1: Fixed Income Fundamentals.
  \item Round 2: NLP and Machine Learning.
\end{itemize}

\textbf{Track B: Investment Analyst Intern (Endowment-Style Role)}
\begin{itemize}
  \item Preparation focus: portfolio construction, valuation, macro sensitivity, and communication.
\end{itemize}

\textbf{Track C: Quantitative Advisory / Consulting Internship (Time-Series Validation)}
\begin{itemize}
  \item Theme mentioned: time-series validation during COVID regime.
\end{itemize}

\textbf{Track D: Portfolio Management / Research Internship (Logistic Regression)}
\begin{itemize}
  \item Questions encountered: logistic regression, L1/L2 regularization, LASSO, and other regularization methods.
\end{itemize}

\section{Behavioral and Fit Preparation}

\RoleProfile{Trading and Commodities Roles}{Commodity merchant, energy trading firm, or commodities desk}

\prompt{1) ``What experience do you have with trading?''}
\begin{itemize}
  \item Use a 60-second narrative: motivation $\to$ what you actually did $\to$ what you learned.
  \item Include at least one concrete example: a backtest, a simulated strategy, paper trading, or risk analysis.
  \item Close with relevance: ``This is why I want a trading analyst role and why I will be effective quickly.''
\end{itemize}

\prompt{2) ``Which commodity are you interested in?''}
\begin{itemize}
  \item Pick one and give a structured reason:
  \begin{enumerate}
    \item Market structure (seasonality, supply chain constraints, multiple drivers).
    \item Data availability (public reports, inventories, weather).
    \item What signals you would explore and why they might exist.
  \end{enumerate}
  \item Avoid vague answers like ``it is interesting''; state 2--3 drivers and a plausible mechanism.
\end{itemize}

\prompt{3) ``What data have you encountered?''}
\begin{itemize}
  \item Answer with type + scale + hard part + how you fixed it.
  \item Examples: time series (prices, volumes), alternative data (text/news), panel data, event data.
  \item Mention one reliability practice: unit tests, data validation checks, reproducible pipeline.
\end{itemize}

\prompt{In-person superday and networking lunch}
\textbf{How to prepare}
\begin{itemize}
  \item Have 2--3 short stories ready: teamwork, disagreement, fast learning, handling ambiguity.
  \item During lunch: be professional, curious, and consistent; avoid controversial topics.
  \item Ask good questions: how they evaluate signals, what makes interns successful, what the team values.
\end{itemize}

\RoleProfile{Technology Audit Summer Analyst (Internal Audit)}{Investment bank or diversified financial institution}

\prompt{1. Internal Audit motivation (60--90 seconds)}
\begin{itemize}
  \item Start with a one-sentence thesis: ``I like work that combines systems thinking, controls, and measurable impact.''
  \item Name 2--3 audit themes: risk identification, control design/testing, stakeholder communication, continuous monitoring.
  \item Give one concrete example from a project where you found/mitigated risk, improved data quality, or built monitoring.
  \item Close with why audit is a fit: breadth across businesses, disciplined process, and learning curve.
\end{itemize}

\prompt{2. Program $\to$ long-term goals}
\begin{itemize}
  \item Goal: become strong in risk/control frameworks and technical auditing (data, systems, process).
  \item Program contributions: rotation exposure, mentorship, real deliverables, cross-functional partners.
  \item Anchor with a specific growth plan: ``I want to strengthen X (controls testing, data lineage, analytics) by doing Y (audits, dashboards, automation).''
\end{itemize}

\prompt{3. GenAI / analytics tools question}
\begin{itemize}
  \item Inventory your tools: 1--2 lines each for Python/SQL/BI tools and what you used them for.
  \item Impact: quantify: reduced manual checks, faster reconciliation, automated anomaly detection.
  \item Future impact: continuous auditing, smarter sampling, narrative reporting, but note governance: model risk, privacy, access control, prompt/data leakage, reproducibility.
  \item Safeguards: human-in-the-loop, audit trails, versioned code, documented assumptions.
\end{itemize}

\RoleProfile{Investment Management Division (Portfolio Solutions Group)}{Asset management firm or institutional investment division}

\prompt{1. Why investment management}
\begin{itemize}
  \item Start with interest: decision-making under uncertainty, long-horizon thinking, portfolio construction.
  \item Show you understand the job: research, risk budgeting, implementation, client constraints.
  \item Provide one example: building a strategy, evaluating factors, measuring risk/return.
\end{itemize}

\prompt{2. Which group and why}
\begin{itemize}
  \item Mention 1--2 groups and tie to specific work outputs (asset allocation, factor exposure control, client portfolio solutions).
  \item Explain your comparative advantage: data, modeling, communication with stakeholders.
  \item Ask/answer: ``What problems does this group solve? What metrics define success?''
\end{itemize}

\prompt{3. Weakness + remediation}
\begin{itemize}
  \item Choose a real but non-fatal weakness (e.g., macro intuition, accounting details, derivatives mechanics).
  \item Show a plan: textbooks/courses, structured reading, building a project/backtest, discussing with peers.
  \item End with outcome: what you can do now that you could not before.
\end{itemize}

\prompt{4. Ethics / ``doing the right thing''}
\begin{itemize}
  \item Use STAR: Situation, Task, Action, Result.
  \item Emphasize escalation path: clarify facts, consult policy/mentor, document, escalate appropriately.
  \item Show balanced judgment: protect clients/firm, avoid assumptions, resolve respectfully.
\end{itemize}

\RoleProfile{Trading Risk Management Intern}{Investment bank, market-making firm, hedge fund, or asset manager}

\prompt{1) ``How do you deal with slackers in a group project?''}
\begin{itemize}
  \item Start with diagnosis: confirm expectations, responsibilities, and constraints (sometimes the issue is unclear scope).
  \item Private alignment first: speak 1:1 to understand blockers; avoid public blaming.
  \item Make work visible: break tasks into small deliverables with dates; use a shared tracker.
  \item Rebalance if needed: reassign tasks to match strengths; define a minimal viable contribution.
  \item Escalate appropriately: if repeated non-delivery, escalate to the project lead with facts and proposed solution.
  \item Close with outcome: emphasize learning and maintaining team morale while protecting deadlines.
\end{itemize}

\prompt{2) ``Why did you pick your master's program?''}
\begin{itemize}
  \item Use a 3-part structure: goal (skills you wanted), fit (courses/projects), evidence (what you produced).
  \item Tie to the role: time series, risk, derivatives, data pipelines, or modeling.
\end{itemize}

\prompt{3) ``Why this role / why this firm?''}
\begin{itemize}
  \item Show understanding of the day-to-day: risk measurement, limit monitoring, stress/scenario work, desk communication.
  \item Mention 1--2 specific learning goals (rates products, swaps, curve risk, model governance).
  \item Connect to past experience: similar risk project, fixed income coursework, or data/automation work.
\end{itemize}

\prompt{4) Resume questions}
\begin{itemize}
  \item Prepare 2--3 ``deep dives'' you can explain: problem, data, method, validation, impact.
  \item Anticipate follow-ups: assumptions, failure cases, and what you would improve.
\end{itemize}

\RoleProfile{Quantitative Technology Intern: Motivation and Prior Experience}{Cross-industry quantitative-finance employers}
\textbf{Questions encountered}
\begin{itemize}
  \item Why are you interested in this position?
  \item Prior banking experience (discussion and impact).
\end{itemize}

\textbf{How to answer (recommended structure)}
\begin{itemize}
  \item Use a 60--90 second narrative: why the role, why the team, and how your skills map to deliverables.
  \item Prior experience: pick one project and give Problem $\to$ Method $\to$ Result (quantified) $\to$ Lesson.
\end{itemize}

\section{Case Studies, Estimation, and Market Intuition}

\RoleProfile{Case Study: Natural Gas Price Prediction}{Commodity merchant, energy trading firm, or commodities desk}
\textbf{What they are testing}
\begin{itemize}
  \item Whether you can translate a business question into a forecasting/feature problem.
  \item Whether you can reason about drivers (fundamentals) and also validate with data.
  \item Whether you avoid common time-series mistakes (leakage, look-ahead bias, wrong validation split).
  \item Whether you can communicate trade-offs: interpretability vs accuracy, simplicity vs complexity.
\end{itemize}

\textbf{A strong answer structure (use this template)}
\begin{enumerate}
  \item \textbf{Clarify target and horizon.} Define what ``price next year'' means: average annual price, year-end spot, front-month futures in 12 months, etc. Confirm frequency (daily/weekly/monthly) and whether predicting level or return.

  \item \textbf{Build a fundamental map of drivers.} Before touching code, list plausible driver categories:
  \begin{itemize}
    \item Supply: production, rig count, pipeline constraints, LNG exports.
    \item Demand: weather (HDD/CDD), industrial usage, power burn.
    \item Storage/inventory: EIA storage levels, seasonal patterns.
    \item Macro/energy complex: crude prices, coal substitution, rate environment.
  \end{itemize}
  Then match dataset columns to these categories and explain why each may lead/lag price.

  \item \textbf{Data checks and feature engineering.} Explain your steps:
  \begin{itemize}
    \item Missingness/outliers: check missing patterns; winsorize extreme errors; do not impute future information.
    \item Stationarity: consider modeling returns/log changes; seasonality features (month-of-year).
    \item Lag features: use only information available at time $t$ to predict $t + \Delta$.
    \item Leakage audit: confirm no post-horizon features sneak in (e.g., future realized inventory).
  \end{itemize}

  \item \textbf{Baseline model first.} Start with a simple benchmark:
  \begin{itemize}
    \item Naive: last value / seasonal average.
    \item Linear regression with lags and seasonal dummies.
    \item Ridge/LASSO for feature selection.
  \end{itemize}
  Explain that beating the baseline is the key, and complexity comes later.

  \item \textbf{Validation (time series).} Use rolling/expanding window:
  \begin{itemize}
    \item Train on early period $\to$ validate on later period, roll forward.
    \item Evaluate MAE/MAPE for levels; directional accuracy if relevant.
    \item Report stability across regimes (high vol vs low vol periods).
  \end{itemize}

  \item \textbf{Interpret results.} Use:
  \begin{itemize}
    \item Feature importance (linear coefficients, permutation importance, SHAP).
    \item Sensitivity analysis: how prediction changes if storage increases, temperature shocks, etc.
    \item Error analysis by season/regime.
  \end{itemize}

  \item \textbf{Deliver the trading interpretation.} Forecast $\to$ how to trade or risk-manage: direction, conviction, risk limits, scenario ranges, and what would falsify the signal.
\end{enumerate}

\textbf{Preparation checklist for this case study}
\begin{itemize}
  \item Understand basic natural gas market drivers (storage seasonality, weather, LNG export sensitivity).
  \item Be able to describe a time-series CV approach without leakage.
  \item Practice: build a small forecasting notebook on any commodity time series with lag features and rolling validation.
  \item Prepare ``model upgrade'' options: tree models, regime features, Bayesian shrinkage; justify only if baseline is solid.
\end{itemize}

\RoleProfile{Brain Teasers / Estimation Questions}{Cross-industry quantitative-finance employers}
\playbooknote{Practice the same topic at two levels: role-specific application questions and general mathematical-foundation questions.}
\textbf{Common prompts}
\begin{itemize}
  \item ``Estimate how many people are flying on airplanes right now in the world.''
  \item ``How many tennis balls fill this room?''
\end{itemize}

\textbf{What they are testing}
Structure, assumptions, arithmetic discipline, and ability to sanity-check.

\textbf{A reusable framework (Fermi template)}
\begin{enumerate}
  \item Define what exactly is being counted (snapshot vs per day, global vs region).
  \item Break the problem into multiplicative components.
  \item Choose reasonable assumptions (say them out loud).
  \item Compute with clean units and intermediate checks.
  \item Provide a range and sensitivity: ``If flights/day is 80k vs 120k, the answer scales linearly.''
\end{enumerate}

\textbf{Example: people flying right now (outline only)}
Flights per day $\times$ average flight duration / 24 hours $\times$ passengers per flight. State assumptions, compute a range, compare to known rough references (global daily passengers).

\RoleProfile{Market Intuition: ``A refinery exploded. What happens?''}{Commodity merchant, energy trading firm, or commodities desk}
\textbf{How to answer}
\begin{enumerate}
  \item Identify the product impacted (gasoline, diesel, jet fuel) and the regional market.
  \item Translate to supply shock: immediate reduction in refinery throughput.
  \item Discuss spreads: crack spreads widen; spot differentials jump; inventory draws.
  \item Substitution and logistics: imports, pipeline reroutes, blending changes.
  \item Time scale: immediate spike, then mean reversion as supply adjusts (unless outage is long).
\end{enumerate}

\section{Coding, Data, and Modeling}

\RoleProfile{Python / pandas Expectations}{Cross-industry quantitative-finance employers}
\begin{itemize}
  \item Data loading, cleaning, joins/merge, \texttt{groupby}, rolling windows, resampling.
  \item Avoid slow loops; use vectorization; explain time complexity if asked.
\end{itemize}

\RoleProfile{Stochastic Process Simulation (How to Present)}{Cross-industry quantitative-finance employers}
\playbooknote{Practice the same topic at two levels: role-specific application questions and general mathematical-foundation questions.}
\begin{enumerate}
  \item Define the model first: for example GBM, OU, AR(1), or GARCH, and say clearly whether it is continuous-time or discrete-time.
  \item State the model-specific evolution rule rather than forcing every process into one generic formula. For example, GBM is commonly written as $S_{t+\Delta} = S_t + \mu S_t\Delta + \sigma S_t\sqrt{\Delta}\,Z$, while OU uses mean reversion $X_{t+\Delta} = X_t + \kappa(\theta - X_t)\Delta + \sigma\sqrt{\Delta}\,Z$.
  \item Generate paths, compute summary statistics, and validate them against what the model implies (mean, variance, stationarity, or reversion speed as appropriate).
\end{enumerate}
\playbooknote{A good interview answer sounds model-aware. Saying ``I would discretize the process'' is not enough; name the drift, diffusion, or discrete update that belongs to the specific process you chose.}

\RoleProfile{Monte Carlo Solution Template}{Cross-industry quantitative-finance employers}
\playbooknote{Practice the same topic at two levels: role-specific application questions and general mathematical-foundation questions.}
\begin{enumerate}
  \item Define random variables and payoff metric.
  \item Simulate $N$ paths; compute estimator $\hat{\theta}$.
  \item Report standard error $SE_c \approx s/\sqrt{N}$ and confidence interval.
  \item Add variance reduction if time allows (antithetic/control variates).
\end{enumerate}

\RoleProfile{SQL Practice Targets (medium-difficulty practice problems)}{Cross-industry quantitative-finance employers}
\begin{itemize}
  \item Joins, window functions (\texttt{ROW NUMBER}, \texttt{LAG}, \texttt{SUM OVER}), grouped aggregations, conditional logic.
  \item Practice writing queries with clear logic and minimal subqueries.
\end{itemize}

\RoleProfile{Missing Data or Alternative Data: What to Do?}{Cross-industry quantitative-finance employers}
\begin{enumerate}
  \item Diagnose missingness (MCAR/MAR/MNAR) and create missing indicators.
  \item Choose imputation method appropriate to data type and leakage risk: median/quantile, forward-fill with caps (time series), KNN/MICE, model-native missing handling (trees).
  \item Validate inside CV folds; stress test sensitivity to imputation choices.
  \item For alternative data: document provenance, stability, and potential bias; monitor drift and revisions.
\end{enumerate}

\RoleProfile{NLP and Machine Learning Interview Questions}{Cross-industry quantitative-finance employers}
\textbf{Questions encountered}
\begin{itemize}
  \item How do you deal with missing data in your project? (median imputation)
  \item Why not use the mean/average?
  \item Label encoding vs one-hot encoding: what is the fundamental difference and how do you apply them?
  \item Walk me through your ML experience on your resume.
  \item How did you do model validation?
\end{itemize}

\textbf{How to answer (step-by-step)}
\begin{enumerate}
  \item \textbf{Missing data.}
  \begin{itemize}
    \item First classify missingness: MCAR/MAR/MNAR; add missingness indicators if predictive.
    \item Median vs mean: median is robust to outliers/skew; mean can be distorted by heavy tails.
    \item Do imputation within CV folds to prevent leakage.
  \end{itemize}

  \item \textbf{Encoding.}
  \begin{itemize}
    \item Label encoding assigns integers and may impose fake ordering; OK for ordinal variables and many tree models.
    \item One-hot avoids ordering, better for linear models/NNs; costs dimensionality and sparsity.
    \item High-cardinality categories: target encoding with K-fold regularization or embeddings.
  \end{itemize}

  \item \textbf{Model validation (generic).}
  \begin{itemize}
    \item Split design: IID vs time series (rolling/expanding windows for time series).
    \item Metrics: choose aligned with task (AUC/logloss for classification; RMSE/MAE for regression).
    \item Calibration: reliability plots; threshold selection; error analysis by segment.
    \item Robustness: stress periods, ablations, and sensitivity checks.
  \end{itemize}

  \item \textbf{Resume ML walkthrough.}
  \begin{itemize}
    \item Problem statement $\to$ data $\to$ feature engineering $\to$ model choice $\to$ validation $\to$ impact.
    \item Mention at least one engineering element: pipeline, reproducibility, tests, monitoring.
  \end{itemize}
\end{enumerate}

\RoleProfile{Time-Series Validation During COVID Regime}{Cross-industry quantitative-finance employers}
\textbf{How to answer (time-series CV best practices)}
\begin{enumerate}
  \item Use rolling/expanding-window splits; never shuffle time.
  \item Prevent leakage for overlapping labels: purged splits + embargo.
  \item Evaluate performance by regime: pre-COVID vs COVID vs post-COVID; include stability metrics.
  \item Document assumptions and add monitoring for drift (PSI, residual tests, volatility changes).
\end{enumerate}

\RoleProfile{Portfolio Management / Research Internship: Logistic Regression}{Asset manager or institutional investment firm}
\textbf{Questions encountered}
\begin{itemize}
  \item How to build logistic regression: inputs, outputs, difficulties.
  \item Did you use L1 and L2? What is LASSO?
  \item Did you use other regularization methods?
\end{itemize}

\textbf{How to answer (step-by-step)}
\begin{enumerate}
  \item \textbf{Model definition.} Binary target $y \in \{0,1\}$; features $x \in \mathbb{R}^d$.
  \[
    p(y = 1 \mid x) = \sigma(\beta_0 + \beta^\top x),
    \qquad
    \sigma(u) = \frac{1}{1 + e^{-u}}.
  \]

  \item \textbf{Training objective.} Negative log-likelihood (cross-entropy):
  \[
    \mathcal{L}(\beta) = \sum_i \left[\log(1 + e^{\beta^\top x_i}) - y_i(\beta^\top x_i)\right].
  \]

  \item \textbf{Difficulties.}
  \begin{itemize}
    \item Class imbalance: use class weights, threshold tuning, PR-AUC.
    \item Multicollinearity: regularize; feature scaling.
    \item Complete separation: L2/Bayesian priors.
    \item Non-stationarity: rolling re-fit, regime validation.
  \end{itemize}

  \item \textbf{Regularization.}
  \begin{itemize}
    \item L2 (Ridge): $\mathcal{L} + \lambda \lVert \beta \rVert_2^2$; stabilizes estimates.
    \item L1 (LASSO): $\mathcal{L} + \lambda \lVert \beta \rVert_1$; induces sparsity.
    \item Elastic Net: combination for correlated features.
  \end{itemize}

  \item \textbf{Validation.} Nested CV; calibration; out-of-time tests; ablation and sensitivity.
\end{enumerate}

\section{Risk, Fixed Income, Options, and Valuation}
\playbooknote{Practice the same topic at two levels: role-specific application questions and general mathematical-foundation questions.}

\RoleProfile{Value at Risk and Risk Modeling}{Investment bank, market-making firm, hedge fund, or asset manager}
\textbf{VaR: a high-quality answer in 60--90 seconds}
\begin{enumerate}
  \item Define VaR precisely (confidence $\alpha$, horizon $H$, loss quantile).
  \item Give 3 methods: historical, parametric (delta-normal), Monte Carlo.
  \item Mention backtesting and limitations (tail risk, assumptions, liquidity).
\end{enumerate}

\textbf{Risk modeling: what to propose}
\begin{itemize}
  \item Factor model (PCA/fundamental factors) for covariance stability.
  \item Heavy tails (Student-$t$), volatility clustering (GARCH/EWMA).
  \item Scenario analysis and stress testing; explain model governance and monitoring.
\end{itemize}

\RoleProfile{Stress Testing Using Data Methods}{Investment bank, market-making firm, hedge fund, or asset manager}
\begin{enumerate}
  \item Choose scenarios: historical replay + hypothetical shocks.
  \item Map shocks to risk factors (rates, spreads, vol, FX) and revalue positions.
  \item Add liquidity and closeout assumptions.
  \item Report: loss distribution, top contributors, and limit breaches.
\end{enumerate}

\RoleProfile{Delta--Gamma Questions}{Investment bank, market-making firm, hedge fund, or asset manager}
\begin{itemize}
  \item Use the second-order approximation:
  \[
    \Delta V \approx \Delta S \cdot \Delta + \frac{1}{2}(\Delta S)^2\Gamma.
  \]
  \item Explain: delta hedging removes first-order risk; gamma remains and drives curvature.
  \item Discuss gamma hedging with an additional option; mention vega/theta trade-offs.
\end{itemize}

\RoleProfile{Trading Risk Management: Technical Questions and How to Answer}{Investment bank, market-making firm, hedge fund, or asset manager}

\prompt{1) Define Value at Risk (VaR).}
\begin{enumerate}
  \item Specify horizon $H$ and confidence $\alpha$.
  \item Define VaR as a quantile of the loss distribution:
  \[
    \mathrm{VaR}_{\alpha}(H) = \inf\{\ell : \mathbb{P}(\mathrm{Loss}_H \le \ell) \ge \alpha\}.
  \]
  \item Mention 3 standard computation methods: historical simulation; parametric (delta-normal); Monte Carlo (full revaluation).
  \item Limitations (one sentence): VaR does not describe tail severity; ES is more informative for tails.
\end{enumerate}

\prompt{2) Pricing mechanism of swaps (explain).}
\begin{enumerate}
  \item A vanilla interest rate swap exchanges fixed vs floating cash flows on a notional $N$.
  \item The par swap rate $S$ is chosen so that the swap has zero value at inception:
  \[
    \mathrm{PV}(\text{fixed leg}) = \mathrm{PV}(\text{float leg}).
  \]
  \item With discount factors $D(t_i)$ and accrual factors $\Delta_i$,
  \[
    \mathrm{PV}(\text{fixed}) = NS\sum_{i=1}^n \Delta_i D(t_i).
  \]
  For a standard par swap, the floating leg PV at inception is approximately
  \[
    \mathrm{PV}(\text{float}) \approx N\bigl(1 - D(t_n)\bigr),
  \]
  hence
  \[
    S = \frac{1 - D(t_n)}{\sum_{i=1}^n \Delta_i D(t_i)}.
  \]
  \item Mention curve bootstrapping (OIS for discounting, IBOR or SOFR-based curves for projection) if asked.
\end{enumerate}

\prompt{3) Compute the present value of a bond on the spot (mental math).}
\begin{enumerate}
  \item Write cash flows: coupons $C$ at $t_1, \ldots, t_n$ plus principal $F$ at $T$.
  \item Discount using yield/spot curve:
  \[
    P = \sum_{i=1}^n C\,D(t_i) + F\,D(T),
    \qquad
    D(t) = e^{-z(t)t}.
  \]
  \item Mental math tips: approximate discount factors for small rates, $e^{-rt} \approx 1 - rt$; group cash flows; check price bounds: if coupon $>$ yield, price $>$ par; if coupon $<$ yield, price $<$ par.
\end{enumerate}

\prompt{4) ``Why are interest rates more volatile as time passes?''}
\begin{itemize}
  \item Clarify which quantity: short rate volatility vs long-maturity yield volatility vs bond price volatility.
  \item A standard interpretation: as horizon increases, uncertainty accumulates. For a random walk-like factor, variance grows with time (e.g., $\mathrm{Var}(W_t) = t$), so longer horizons can exhibit larger uncertainty in integrated quantities.
  \item For bonds, price sensitivity to rate changes is driven by duration: longer maturity generally has higher duration, hence larger price moves for a given yield shock.
  \item Long yields can be smoother if they average expected future short rates; in practice, term premium and regime changes still create meaningful long-horizon variance.
\end{itemize}

\prompt{5) Statistics question about distributions (generic approach).}
\begin{itemize}
  \item Identify distribution family and parameters; compute mean/variance; translate the question into a probability or expectation.
  \item Use conditioning, CLT, or tail approximations if exact forms are hard.
  \item Provide sanity checks (limits, symmetry, scaling).
\end{itemize}

\RoleProfile{Fixed Income Fundamentals}{Investment bank, asset manager, insurer, or institutional investor}

\prompt{1. What is a bond?}
\begin{itemize}
  \item Definition: a security that pays scheduled cash flows (coupons + principal).
  \item Pricing: present value of cash flows discounted by the term structure.
  \item Risks: rate risk, credit risk, liquidity, embedded options (call/prepay).
\end{itemize}

\prompt{2. What is duration?}
\begin{itemize}
  \item Macaulay duration = PV-weighted average time of cash flows.
  \item Modified duration approximates first-order price sensitivity: $\Delta P/P \approx -D_{\mathrm{mod}}\,\Delta y$.
  \item Mention convexity as second-order correction.
\end{itemize}

\prompt{3. How to calculate duration (step-by-step).}
\begin{enumerate}
  \item Build cash flow schedule $\{(t_i, CF_i)\}$.
  \item Discount each cash flow using spot rates $z(t_i)$ or a yield $y$: $PV_i = CF_i e^{-z(t_i)t_i}$.
  \item Compute price $P = \sum_i PV_i$.
  \item Macaulay duration $D_{\mathrm{Mac}} = \dfrac{1}{P}\sum_i t_i PV_i$.
  \item Modified duration (consistent compounding convention): $D_{\mathrm{mod}} = D_{\mathrm{Mac}}/(1 + y)$.
  \item For curve risk: compute key-rate durations by bumping one tenor at a time.
\end{enumerate}

\prompt{4. Yield curve (what to say when asked ``tell me more'').}
\begin{itemize}
  \item Define: mapping maturity $t$ to spot/zero rates $z(t)$ or discount factors $D(t)$.
  \item Construction: bootstrap from liquid instruments (bills/OIS/swaps); interpolate smoothly.
  \item Risk: interpret moves via level/slope/curvature; manage via key-rate DV01.
  \item Use cases: pricing, hedging, forward rates, scenario/stress testing.
\end{itemize}

\RoleProfile{Investment Analyst Intern (Endowment-Style Role)}{Endowment, foundation, outsourced investment office, or institutional investor}
\textbf{Preparation focus (typical)}
\begin{itemize}
  \item Portfolio construction: diversification, factor exposures, drawdown and risk budgeting.
  \item Valuation: basic multiples, cash flow logic, understanding macro sensitivity.
  \item Communication: ability to summarize an investment thesis concisely and defend assumptions.
\end{itemize}

\RoleProfile{Valuation and Expected Value Problems (EV/EBITDA + Default Risk)}{Investment bank, asset manager, insurer, or institutional investor}
\textbf{How to answer these questions (step-by-step)}
\begin{enumerate}
  \item Define the payoff clearly. Identify states of the world (default / no default; price up/down), and write payoff $X$ in each state.
  \item Compute fair price as discounted expectation. For horizon $T$ with discount rate $r$:
  \[
    P_0 = e^{-rT}\,\mathbb{E}[X] = e^{-rT}\sum_s p_s X_s.
  \]
  \item Include recovery and loss-given-default when needed. If default probability is $p$ and recovery is $R$ (fraction of face value $F$):
  \[
    \mathbb{E}[X] = (1 - p)\cdot F + p\cdot (RF).
  \]
  \item Connect EV/EBITDA to enterprise value. Recall definitions:
  \[
    EV = \text{Market Cap} + \text{Debt} - \text{Cash},
  \]
  and the multiple EV/EBITDA is used to infer EV from EBITDA:
  \[
    EV = (EV/EBITDA)_{\mathrm{peer}} \times EBITDA.
  \]
  Then back out equity value via $\text{Equity} = EV - \text{Debt} + \text{Cash}$.
  \item Sanity checks. Make sure units match, probabilities sum to 1, and price lies between extreme outcomes (after discounting).
\end{enumerate}

\section{Assessment Formats and Strategy}

\RoleProfile{Preparation Focus Across Mentioned Assessment Formats}{Multiple quantitative-finance institution types}
\begin{itemize}
  \item Time pressure: practice doing clean reasoning quickly and writing concise solutions.
  \item Coding: Python data structures, complexity, vectorization, and debugging speed.
  \item Math: probability distributions, calculus (derivatives/integrals), linear algebra basics.
  \item Communication: in group exercises, narrate assumptions, assign roles, and keep the team aligned.
\end{itemize}

\RoleProfile{Proctored Math Online Assessment (Quant Research Intern)}{Hedge fund, quantitative investment manager, or proprietary trading firm}
\playbooknote{For timed breadth, draw heavily from a suitable practice set: calculus/linear algebra, probability, Markov chains, martingales, and Brownian motion. Keep the easier questions to a 3--5 minute target and reserve longer dynamic-programming questions for the end.}
\textbf{What this format is testing}
\begin{itemize}
  \item Speed and accuracy under time pressure.
  \item Breadth across probability/combinatorics (not just ``light'' mental math).
  \item Ability to pick the right technique quickly (symmetry, conditioning, recursion, linearity of expectation).
\end{itemize}

\textbf{How to approach each question (repeatable template)}
\begin{enumerate}
  \item Classify the question in 5--10 seconds: EV? counting? conditional probability? recursion? geometry?
  \item Write the key random variable(s). State exactly what you are computing (e.g., $\mathbb{E}[X]$, $\mathbb{P}(A \mid B)$).
  \item Look for shortcuts first.
  \begin{itemize}
    \item Linearity of expectation even without independence.
    \item Symmetry (exchangeable outcomes).
    \item Complement (compute $1 - \mathbb{P}(\cdot)$ if easier).
    \item Condition on the first step for processes.
    \item Indicator variables for counts.
  \end{itemize}
  \item Compute cleanly. Keep fractions exact; postpone decimals.
  \item Sanity-check. Is the answer in $[0,1]$? Does it match edge cases?
\end{enumerate}

\textbf{Preparation plan (90 minutes / 16 questions)}
\begin{itemize}
  \item Target pacing: $\sim$5 minutes/question average; plan 10--15 minutes buffer for 2--3 hard ones.
  \item Drill sets: binomial/negative binomial, geometric, hypergeometric; coupon collector variants; Markov chains with 2--3 states; classical combinatorics.
  \item Brain teasers: practice Fermi estimation, invariants, and ``constructive'' logic.
  \item Build an error log of common misses: off-by-one, conditional sample space, double counting.
\end{itemize}

\RoleProfile{24-Hour Take-Home Coding Assessment (Quant Research Intern)}{Hedge fund, quantitative investment manager, or proprietary trading firm}
\textbf{How to structure your solution (what reviewers look for)}
\begin{enumerate}
  \item Restate requirements. Write a brief spec: inputs, outputs, constraints, edge cases.
  \item Design first. Identify the core abstractions and data flow.
  \begin{itemize}
    \item Quote (bid, ask, timestamp), Order, Trade, Venue, Book.
    \item A pricing/arb module: detect violations like bid $>$ ask across venues or implied cross-rate arbitrage.
  \end{itemize}
  \item Implement incrementally. Start with minimal features + tests; then add enhancements.
  \item Correctness first. Unit tests for each module; golden examples.
  \item Performance and clarity. Avoid premature micro-optimizations, but use appropriate data structures (dict/heap/Deque).
  \item Documentation. README: how to run, assumptions, complexity.
\end{enumerate}

\textbf{Spreads and arbitrage: what to explain clearly}
\begin{itemize}
  \item Spread = ask -- bid; midpoint = (ask + bid)/2; crossing markets if bid $\ge$ ask.
  \item Cross-venue arb: if venue A bid $>$ venue B ask, buy at B ask, sell at A bid (after costs/latency).
  \item Include frictions: fees, slippage, minimum trade size, stale quotes; otherwise ``arb'' is illusory.
\end{itemize}

\textbf{Take-home time management}
\begin{itemize}
  \item First 1 hour: plan + write skeleton + tests.
  \item Next 3--6 hours: implement core logic.
  \item Last 1--2 hours: refactor, add edge cases, improve README, run lint/tests.
\end{itemize}

\RoleProfile{Proctored Python Coding (Mostly pandas)}{Cross-industry quantitative-finance employers}
\textbf{How to answer efficiently (pandas assessment playbook)}
\begin{enumerate}
  \item Understand the schema. \texttt{df.info()}, \texttt{df.head()}, missingness and unique counts.
  \item Clean deterministically.
  \begin{itemize}
    \item Parse dates with \texttt{pd.to\_datetime}.
    \item Enforce dtypes: \texttt{astype('int64')} / nullable \texttt{'Int64'} when missing values exist.
    \item Handle missing values explicitly and consistently.
  \end{itemize}
  \item Wrangling patterns to know.
  \begin{itemize}
    \item \texttt{merge/join}, \texttt{groupby().agg()}, \texttt{pivot\_table}, \texttt{sort\_values}.
    \item Window operations: \texttt{rolling}, \texttt{expanding}, \texttt{shift}, \texttt{diff}.
  \end{itemize}
  \item Baseline model first. Train/test split (or time-based split), then logistic/linear model baseline.
  \item Evaluation. Choose metric aligned with task: RMSE/MAE for regression; AUC/F1/logloss for classification.
  \item Edge-case checks. Duplicate keys, leakage, label imbalance, unexpected categories.
\end{enumerate}

\textbf{Common pitfalls to avoid}
\begin{itemize}
  \item Confusing \texttt{int64} vs pandas nullable \texttt{Int64}; failing type checks in hidden tests.
  \item Silent merges that increase row count due to many-to-many keys.
  \item Leakage (using target-derived features or future information).
  \item Using chained assignment or not handling NaNs before casting.
\end{itemize}

\textbf{A minimal ``dtype-safe'' template (snippet)}
\begin{codeblock}
import pandas as pd

def coerce_int64(series: pd.Series) -> pd.Series:
    # Use pandas nullable integer if NaNs exist
    if series.isna().any():
        return series.astype("Int64")
    return series.astype("int64")
\end{codeblock}

\RoleProfile{Long-Form Game-Theory Puzzle (Card Game Inference)}{Proprietary trading or electronic market-making firm}
\textbf{What this format is testing}
\begin{itemize}
  \item Strategic reasoning under uncertainty (game-theoretic optimality).
  \item Bayesian-style inference: actions provide information about hidden states.
  \item Ability to write detailed justification: not just the answer, but why it follows from optimal play.
  \item Consistency: your beliefs and your recommended strategy should align.
\end{itemize}

\textbf{How to solve these problems (recommended template)}
\begin{enumerate}
  \item Restate the game formally.
  \begin{itemize}
    \item Define the hidden state (the card type or hand).
    \item Define the action set (what a player can do at each step).
    \item Define payoffs and stopping conditions.
    \item Identify what information each player has (information sets).
  \end{itemize}
  \item Write down the optimality principle.
  \begin{itemize}
    \item If the other player is rational, their action should maximize expected payoff given their information.
    \item Use backward induction if the game has a finite horizon; otherwise use threshold logic or value functions.
  \end{itemize}
  \item Connect actions to beliefs (inference).
  \begin{itemize}
    \item For each possible card/hand, compute what the rational action would be.
    \item Eliminate states inconsistent with the observed action.
    \item If multiple states remain, assign probabilities using Bayes' rule:
    \[
      \mathbb{P}(H \mid A) = \frac{\mathbb{P}(A \mid H)\mathbb{P}(H)}{\sum_h \mathbb{P}(A \mid h)\mathbb{P}(h)}.
    \]
  \end{itemize}
  \item Compute expected values and compare actions.
  \begin{itemize}
    \item For each action, compute EV under your posterior beliefs.
    \item Pick the action with the highest EV (or highest utility if risk-averse utility is specified).
    \item Provide a sensitivity check: ``If the prior were different, the threshold would move.''
  \end{itemize}
  \item Write the justification like a proof.
  \begin{itemize}
    \item State assumptions (rationality, priors, tie-breaking).
    \item Show each elimination step or EV comparison clearly.
    \item End with a concise conclusion: posterior odds and the recommended action.
  \end{itemize}
\end{enumerate}

\textbf{How to prepare}
\begin{itemize}
  \item Practice expected value games: optimal stopping, mixed strategies, threshold policies.
  \item Practice inference from actions: signaling games, ``rational agent'' puzzles.
  \item Practice writing: for each puzzle, force yourself to write a 5--10 sentence justification.
  \item Time management for long sessions: take structured notes; freeze definitions; avoid re-reading rules repeatedly.
\end{itemize}

\RoleProfile{Cognitive Ability Test (Pattern + Memory Games)}{Proprietary trading or electronic market-making firm}
\textbf{What they are testing}
\begin{itemize}
  \item Speed-accuracy tradeoff under strict time limits.
  \item Working memory capacity and attention control.
  \item Visual reasoning: transformations (rotation, reflection, scaling) and rule extraction.
\end{itemize}

\textbf{How to perform better (practical tips)}
\begin{enumerate}
  \item Pattern tasks: scan for invariants first (counts, symmetry, parity), then transformations (rotate/flip/shift), then arithmetic rules (increment, alternate).
  \item Memory tasks: chunk sequences (e.g., 2-2-3), map to a grid coordinate system, and keep a consistent reading order.
  \item Pacing: do not get stuck. If uncertain, guess quickly and move on; the scoring is often throughput-sensitive.
  \item Practice: short daily drills outperform occasional long sessions.
\end{enumerate}

\RoleProfile{Valuation and Expected Value Problems in Online Assessments}{Investment bank, asset manager, insurer, or institutional investor}
\textbf{How to prepare}
\begin{itemize}
  \item Drill expected value pricing with discrete scenarios (default/recovery, barrier events, partial payouts).
  \item Memorize core valuation identities (EV, equity value bridge).
  \item Practice quick accounting interpretations: what changes EV vs EBITDA, how leverage affects equity.
\end{itemize}

\section{Selected Technical Practice Questions}
\playbooknote{The seven worked examples below stay compact on purpose. After each example, add one calculus, linear-algebra, or probability variation.}

\RoleProfile{Macro Hedge Fund --- Technical Questions (with solutions)}{Macro hedge fund}

\prompt{Q1. Differential equations (basic)}
\textbf{Problem.} Solve the linear ODE $y'(x) = ay(x) + b$ with $y(0) = y_0$.

\textbf{Solution.} Rewrite $y' - ay = b$. Multiply by integrating factor $e^{-ax}$:
\[
  \frac{d}{dx}\bigl(e^{-ax}y\bigr) = be^{-ax}.
\]
Integrate: $e^{-ax}y - y_0 = \dfrac{b}{-a}(e^{-ax} - 1)$ if $a \ne 0$. Thus
\[
  y(x) = \left(y_0 + \frac{b}{a}\right)e^{ax} - \frac{b}{a}
  \qquad (a \ne 0),
  \qquad
  y(x) = y_0 + bx
  \qquad (a = 0).
\]
\playbooknote{This is a standard integrating-factor problem. The safest way to avoid sign mistakes in an interview is to check the initial condition immediately at $x=0$ before you move on.}

\prompt{Q2. Two random points on a circle: chord longer than the radius}
\textbf{Problem.} Pick two points uniformly at random on a circle of radius $R$. What is the probability the chord length exceeds $R$?

\textbf{Solution.} Let $\theta \in [0,\pi]$ be the smaller central angle between the points. For uniform points, $\theta$ is uniform on $[0,\pi]$. Chord length is $L = 2R\sin(\theta/2)$. We want $L > R$:
\[
  2R\sin(\theta/2) > R
  \Longleftrightarrow
  \sin(\theta/2) > \frac12
  \Longleftrightarrow
  \theta/2 > \frac{\pi}{6}
  \Longleftrightarrow
  \theta > \frac{\pi}{3}.
\]
Therefore
\[
  \mathbb{P}(L > R) = \frac{\pi - \pi/3}{\pi} = \frac{2}{3}.
\]

\prompt{Q3. Differentiate an integral with variable limit (Fundamental Theorem)}
\textbf{Problem.} If $F(x) = \int_0^x s\,ds$, compute $F'(x)$.

\textbf{Solution.} By the Fundamental Theorem of Calculus, $F'(x) = x$. (Also $F(x) = \tfrac12 x^2$ so derivative is $x$.) Hence $F'(x) = x$.

\prompt{Q4. Four fair dice: probability the maximum equals 3}
\textbf{Problem.} Roll 4 fair six-sided dice. What is $\mathbb{P}(\max = 3)$?

\textbf{Solution.}
\[
  \mathbb{P}(\max \le 3) = \left(\frac{3}{6}\right)^4 = \left(\frac12\right)^4 = \frac1{16},
  \qquad
  \mathbb{P}(\max \le 2) = \left(\frac{2}{6}\right)^4 = \left(\frac13\right)^4 = \frac1{81}.
\]
Thus
\[
  \mathbb{P}(\max = 3) = \mathbb{P}(\max \le 3) - \mathbb{P}(\max \le 2) = \frac1{16} - \frac1{81} = \frac{65}{1296}.
\]

\prompt{Q5. Orthogonal matrix characterization}
\textbf{Problem.} For a matrix $A$, what does it mean for $A$ to be orthogonal?

\textbf{Solution.} $A$ is orthogonal iff it preserves dot products and lengths:
\[
  A^\top A = I
  \qquad
  (\text{equivalently } A^{-1} = A^\top).
\]
Then $\lVert Ax \rVert_2 = \lVert x \rVert_2$ and $(Ax)^\top(Ay) = x^\top y$ for all $x,y$.
\playbooknote{If the interviewer asks for intuition, say ``orthogonal matrices preserve geometry'' before you say ``$A^\top A = I$''. That usually lands better than giving the matrix identity with no interpretation.}

\prompt{Q6. Newton's method convergence speed (conceptual)}
\textbf{Problem.} Newton iteration for $f(x) = \tanh(x^4 - x^2)$. Which root converges faster and why?

\textbf{Solution.} Newton's method has quadratic local convergence to a simple root where $f'(x^\star) \ne 0$. If a root is multiple or has near-zero derivative, convergence slows. Here $\tanh(u) = 0$ iff $u = 0$, so roots come from $x^4 - x^2 = x^2(x^2 - 1) = 0$:
\[
  x^\star \in \{-1,0,1\}.
\]
Near $x = 0$, $u = x^4 - x^2 \approx -x^2$ is second order in $x$, so $f(x) \approx -x^2$ and $f'(0) = 0$, which behaves like a multiple root and leads to slower convergence. Near $x = \pm 1$, the derivative is nonzero (simple roots), so Newton converges faster there. Thus convergence is faster near $x = \pm 1$ than near $x = 0$.

\prompt{Q7. $X, Y \sim \mathrm{Unif}(0,1)$: compute $\mathbb{P}(X^2 > Y)$}
\textbf{Problem.} If $X,Y$ are independent uniform on $(0,1)$, find $\mathbb{P}(X^2 > Y)$.

\textbf{Solution.} Interpret geometrically on the unit square:
\[
  \mathbb{P}(X^2 > Y) = \int_0^1 \mathbb{P}(Y < x^2)\,dx = \int_0^1 x^2\,dx = \left[\frac{x^3}{3}\right]_0^1 = \frac13.
\]

\section{Preparation Checklists}

\RoleProfile{Case Study Checklist: Natural Gas Forecasting}{Commodity merchant, energy trading firm, or commodities desk}
\begin{itemize}
  \item Understand basic natural gas market drivers (storage seasonality, weather, LNG export sensitivity).
  \item Be able to describe a time-series CV approach without leakage.
  \item Practice: build a small forecasting notebook on any commodity time series with lag features and rolling validation.
  \item Prepare ``model upgrade'' options: tree models, regime features, Bayesian shrinkage; justify only if baseline is solid.
\end{itemize}

\RoleProfile{General Assessment Preparation Checklist}{Multiple quantitative-finance institution types}
\begin{itemize}
  \item Time pressure: practice doing clean reasoning quickly and writing concise solutions.
  \item Coding: Python data structures, complexity, vectorization, and debugging speed.
  \item Math: probability distributions, calculus (derivatives/integrals), linear algebra basics.
  \item Communication: in group exercises, narrate assumptions, assign roles, and keep the team aligned.
\end{itemize}

\RoleProfile{Trading Risk Management Preparation Checklist}{Investment bank, market-making firm, hedge fund, or asset manager}
\begin{itemize}
  \item Rates basics: discount factors, bond PV, duration/convexity, swap par rate.
  \item Risk basics: VaR definitions + methods, stress testing, limit monitoring.
  \item Data: missingness handling, leakage avoidance, and time-series validation.
  \item Behavioral: STAR stories for teamwork conflict, motivation, and project deep dives.
\end{itemize}

\RoleProfile{Investment Analyst / Endowment-Style Preparation Focus}{Endowment, foundation, outsourced investment office, or institutional investor}
\begin{itemize}
  \item Portfolio construction: diversification, factor exposures, drawdown and risk budgeting.
  \item Valuation: basic multiples, cash flow logic, understanding macro sensitivity.
  \item Communication: ability to summarize an investment thesis concisely and defend assumptions.
\end{itemize}


\newpage

\section{Role Maps and Preparation Priorities}
\begin{longtable}{@{}p{0.22\textwidth}p{0.30\textwidth}p{0.22\textwidth}p{0.20\textwidth}@{}}
\toprule
\textbf{Role track} & \textbf{Typical stages} & \textbf{Most tested themes} & \textbf{Prep priority if time is short}\\
\midrule
\endfirsthead
\toprule
\textbf{Role track} & \textbf{Typical stages} & \textbf{Most tested themes} & \textbf{Prep priority if time is short}\\
\midrule
\endhead
\bottomrule
\endfoot
Commodities trading & HR screen, market case, technical round, superday, networking lunch & market intuition, time-series structure, trading motivation, fast arithmetic & build one commodity case, one Fermi framework, and one clean trading story\\
Quant risk analyst & fit + technical, risk modelling, stress testing, option-risk discussion & VaR, factor risk, scenario design, delta-gamma thinking & be able to give a 60-second and a 90-second VaR answer\\
Internal audit / tech audit & pre-recorded video interview, behaviorals, analytics tools question & controls mindset, structured communication, data tools with governance & prepare three STAR stories and one analytics-use case with safeguards\\
Investment management / portfolio solutions & PRVI, group preference, weakness, ethics & portfolio thinking, client constraints, judgment, remediation & have one investment or factor example and one weakness-remediation story\\
Quant research intern & math OA, take-home coding, pandas or modelling test & probability speed, coding completeness, debugging, notation discipline & drill timed probability sets and one pandas notebook end-to-end\\
Trading risk management & HR plus mixed manager rounds & swaps, bonds, rate risk, behavior under ambiguity & review rates basics and rehearse why-this-role answers with desk language\\
Quant technology / advisory / PM research & fixed income fundamentals, ML validation, regularization & bond math, data validation, leakage control, logistic regression intuition & build one short fixed-income explainer and one validation explainer\\
\end{longtable}

\playbooknote{A useful preparation heuristic is to match every role to three buckets: what they can ask anyone, what they ask because of the desk or asset class, and what they ask because of the assessment format. Once those buckets are visible, the interview stops feeling random.}


\section{Preparation Systems}
\RoleProfile{14-day sprint plan}{Multiple quantitative-finance institution types}
\begin{enumerate}
  \item Days 1--2: inventory the target roles, map them to the companion bank, and highlight the 10 prompts or topics most likely to recur.
  \item Days 3--5: build the story bank - motivation, teamwork conflict, ambiguity, weakness, ethics, and one deep technical project.
  \item Days 6--8: run technical short-answer drills on VaR, swaps, bond PV, duration, validation, missing data, and logistic regression; add one 8--10 question ZQ set each day from the matching chapter.
  \item Days 9--10: do one timed math set using QB/ZQ records and one timed coding or pandas set under realistic conditions.
  \item Days 11--12: rehearse one market or case-study answer out loud until it becomes procedural rather than improvised.
  \item Days 13--14: run two mock interviews, tighten weak openings, and prepare one-page cheat sheets for role-specific topics.
\end{enumerate}

\RoleProfile{28-day build plan}{Multiple quantitative-finance institution types}
\begin{enumerate}
  \item Week 1: foundations - use a suitable practice set for brain teasers, calculus, linear algebra, probability, and stochastic-process refresh; build the behavioral story inventory in parallel.
  \item Week 2: technical frameworks - pair the QB risk/derivatives material with a suitable practice set, then add validation and desk-specific concepts.
  \item Week 3: assessments - timed math, pandas, take-home simulation, and group-exercise communication drills; use a suitable practice set for algorithms and numerical methods.
  \item Week 4: consolidation - mocks, role-specific reading, résumé deep dives, and day-before packs for each target firm.
\end{enumerate}

\RoleProfile{Day-before checklist}{Multiple quantitative-finance institution types}
\begin{itemize}
  \item Rehearse your opening for ``tell me about yourself'' and ``why this role'' one last time.
  \item Revisit only short-answer frameworks and formulas; avoid learning brand-new material.
  \item Prepare one example each for teamwork, conflict, fast learning, ethics, and technical depth.
  \item For assessment formats, decide your pacing rule in advance: when to skip, when to guess, and when to return.
\end{itemize}

\RoleProfile{Day-of checklist}{Multiple quantitative-finance institution types}
\begin{itemize}
  \item Start every answer by naming the structure you are using: STAR, define-and-caveat, or clarify-and-decompose.
  \item If you are uncertain, say your assumption cleanly and continue. Ambiguity handled well scores better than silence.
  \item Keep a short list of anchor phrases ready: ``let me define the object first'', ``the main trade-off is'', ``the risk is leakage'', ``the desk implication is''.
\end{itemize}


\section{Story Bank Builder}
\RoleProfile{Core stories to prepare}{Cross-industry quantitative-finance employers}
\begin{longtable}{@{}p{0.24\textwidth}p{0.32\textwidth}p{0.38\textwidth}@{}}
\toprule
\textbf{Story type} & \textbf{What the interviewer is really probing} & \textbf{What your story needs}\\
\midrule
\endfirsthead
\toprule
\textbf{Story type} & \textbf{What the interviewer is really probing} & \textbf{What your story needs}\\
\midrule
\endhead
\bottomrule
\endfoot
Motivation / why this role & seriousness of intent and role understanding & one real reason, one relevant example, one explicit fit statement\\
Conflict / slacker / disagreement & judgment and maturity under friction & diagnosis, private conversation, concrete action, outcome without drama\\
Weakness + remediation & self-awareness and growth rate & a real gap, a believable fix, and evidence that the fix changed your behaviour\\
Ethics / doing the right thing & escalation judgment and integrity & facts first, policy or mentor check, respectful escalation, documented result\\
Technical project deep dive & real ownership and rigor & problem, data, method, validation, impact, and one thing you would improve\\
Ambiguity / fast learning & comfort with incomplete information & how you framed the problem, prioritized, and learned while moving\\
\end{longtable}

\RoleProfile{One-minute answer template}{Cross-industry quantitative-finance employers}
\begin{enumerate}
  \item Opening sentence: the one thing the interviewer should remember.
  \item Context: only the facts needed to understand the example.
  \item Your action: the decision or work you personally owned.
  \item Result: quantified if possible, otherwise specific.
  \item Relevance: why this example matters for the role in front of you.
\end{enumerate}


\section{Mock Interview Packs}
\RoleProfile{Pack A: Trading / commodities}{Commodity merchant, energy trading firm, or commodities desk}
\playbooknote{Select a compact set of matching probability, trading, risk, or coding questions from your practice materials.}
\begin{enumerate}
  \item Why trading, and why this commodity in particular?
  \item Walk me through the natural gas forecasting case from target definition to validation.
  \item Estimate a market size or a real-world quantity with explicit assumptions.
  \item A refinery explodes. What happens in the first hour, the first day, and the first week?
  \item Tell me about a time you used data to make or defend a decision.
\end{enumerate}

\RoleProfile{Pack B: Risk / fixed income}{Investment bank, asset manager, insurer, or institutional investor}
\playbooknote{Select a compact set of matching probability, trading, risk, or coding questions from your practice materials.}
\begin{enumerate}
  \item Define VaR in 60 seconds and tell me one limitation.
  \item Explain how a swap is priced and what the par rate means.
  \item Price a bond on the spot and tell me how coupon vs. yield affects premium or discount.
  \item How would you stress test a portfolio across rates, spreads, vol, and liquidity?
  \item Why this role, and what part of rate or portfolio risk interests you most?
\end{enumerate}

\RoleProfile{Pack C: Data / ML / quant technology}{Cross-industry quantitative-finance employers}
\playbooknote{Select a compact set of matching probability, trading, risk, or coding questions from your practice materials.}
\begin{enumerate}
  \item How did you handle missing data, and why did you choose that method?
  \item Explain label encoding versus one-hot encoding and when each is appropriate.
  \item Walk me through your validation design and how you prevented leakage.
  \item Describe one ML project end-to-end: problem, data, modelling, validation, and impact.
  \item How would you explain logistic regression regularization to a non-specialist stakeholder?
\end{enumerate}


\RoleProfile{Pack D: General Quant Foundations}{Multiple quantitative-finance institution types}
\playbooknote{Use this pack when the role is broad or the interviewer is likely to test mathematical fundamentals independent of a specific desk.}
\begin{enumerate}
  \item One structural brain teaser from a suitable practice set; explain the invariant or backward-induction step before calculating.
  \item One calculus or linear-algebra question from a suitable practice set; derive rather than quote the answer.
  \item Two probability questions from a suitable practice set, with at least one using conditioning or expected value.
  \item One stochastic-process question from a suitable practice set; state the process, stopping rule, or martingale clearly.
  \item One finance question from a suitable practice set or a role-specific derivatives or risk question.
  \item One algorithm/numerical question from a suitable practice set; give correctness and complexity before implementation details.
\end{enumerate}
\section{Conclusion}
Interview preparation becomes more reliable when it is treated as a system rather than a collection of memorized answers. Identify the institution type and role, classify the format, choose a compact answer structure, make assumptions explicit, and finish with an interpretation or decision. Technical depth matters, but so do pacing, error checks, and communication.

\section{Review Questions}
\begin{enumerate}
  \item How should preparation differ among a proprietary trading firm, an investment bank, an asset manager, and a quantitative investment manager?
  \item What is the shortest complete structure for a behavioral answer, and what evidence makes it credible?
  \item How would you design a leakage-safe validation plan for a time-series forecasting case?
  \item What should a candidate say before beginning the arithmetic in an estimation problem?
  \item How do VaR, stress testing, and scenario analysis complement one another?
  \item What changes when a coding task is a timed online assessment rather than a collaborative live interview?
  \item Build a two-week practice schedule that combines behavioral stories, technical drills, cases, and mock interviews.
\end{enumerate}
\chapter{Selected Interview Questions and Worked Solutions}

\section{Selection and Study Method}
This chapter contains 50 representative interview questions: 35 role-linked questions and 15 cross-industry foundation questions. The selection favors recurring concepts, transferable solution methods, and coverage across trading, research, risk, investment, data, and technology roles. Repetitive variants and unusually specialized problems have been omitted.

Specific employer names and identifying details are not shown. For role-linked questions, the institution type, role, and difficulty are retained. Foundation questions are labeled for cross-industry quantitative-finance employers.

Use a three-pass method. On the first pass, answer the problem without notes and explain the setup aloud. On the second pass, compare the method with the worked solution and record the exact failure mode. On the third pass, answer a variation under time pressure. A question is not mastered merely because its solution looks familiar.

\begin{longtable}{@{}L{0.58\textwidth}r@{}}
\toprule
\textbf{Selected area} & \textbf{Questions} \\
\midrule
Probability, puzzles, combinatorics, and decision-making & 8 \\
Statistics, machine learning, linear algebra, and time series & 5 \\
Stochastic calculus, options, and hedging & 5 \\
Risk, credit, fixed income, and portfolio management & 5 \\
Trading, market microstructure, signals, and backtesting & 4 \\
Algorithms, data structures, coding, and systems & 6 \\
Behavioral, communication, and governance & 2 \\
Cross-industry mathematical and finance foundations & 15 \\
\midrule
\textbf{Total} & \textbf{50} \\
\bottomrule
\end{longtable}

\section*{Role-Linked Questions}
\addcontentsline{toc}{section}{Role-Linked Questions}
\section{Probability, Puzzles, and Combinatorics}
\ChapterIntro{Use these questions to practice complete solutions, concise explanations, and likely follow-up questions.}

\subsection{Logic puzzles and counting arguments}

\InterviewQuestionHeader{Question 1}{Options market-making and proprietary trading firm}{Quantitative Trader}{Logic / Measurement Puzzle}{Easy}
\textbf{Problem.} You have a 1-liter bottle and a 1.5-liter bottle. How do you measure out exactly 0.5 liters?

\textbf{Solution.}

Fill the 1.5 L bottle completely, then pour from it into the 1 L bottle until the 1 L bottle is full. The amount left in the 1.5 L bottle is
\[
1.5 - 1.0 = 0.5 \text{ L}.
\]

\begin{EditorialNote}


\textbf{Core skill tested.} This is a structured elimination problem rather than a formula dump. The interviewer is looking for disciplined bookkeeping: define the unknowns, translate each statement into a constraint, and use consistency to rule out impossible cases without losing track of the assumptions.

\textbf{How to present it in interview.} State the unknowns and the rule of the game first, then convert each sentence into a constraint you can reuse. Interviewers like hearing the consistency checks out loud because it shows that you are controlling the search space instead of guessing and backfitting.

\textbf{Common traps.}
\begin{itemize}
\item Treating each sentence in isolation instead of writing a small system of constraints.
\item Forgetting that truth/lie or ordering rules depend on who is speaking to whom.
\item Doing arithmetic too early before the structural consistency conditions are fixed.
\end{itemize}

\textbf{Related follow-ups.}
\begin{itemize}
\item Can you solve the same puzzle by contradiction or by ordering the cases differently?
\item What is the minimum set of constraints that actually determines the unique solution?
\end{itemize}


\end{EditorialNote}
\QuestionSpacer

\InterviewQuestionHeader{Question 2}{Quantitative investment management firm}{Quantitative Researcher}{Puzzle / Tournament Reasoning}{Medium}
\textbf{Problem.} You have 25 horses and can race 5 at a time; there is no timer (only relative ordering within a race). What is the minimum number of races needed to determine the top 3 fastest horses?

\textbf{Solution.}

The minimum number of races is \(\boxed{7}\).

\begin{enumerate}
\item Split the 25 horses into 5 groups of 5, and race each group once. This takes 5 races. Let the finishing orders within the groups be
\[
A_1 > A_2 > A_3 > A_4 > A_5,\quad \ldots,\quad E_1 > E_2 > E_3 > E_4 > E_5.
\]

\item Race the 5 group winners \(A_1,B_1,C_1,D_1,E_1\). Suppose the order is
\[
A_1 > B_1 > C_1 > D_1 > E_1.
\]
This is the 6th race.

From this, every horse in groups \(D\) and \(E\) is eliminated. Also, \(B_3,B_4,B_5\), \(C_2,C_3,C_4,C_5\), and \(A_4,A_5\) cannot be in the top 3 overall.

\item The only remaining candidates for the top 3 are
\[
\{A_1, A_2, A_3, B_1, B_2, C_1\}.
\]
Since \(A_1\) already beat \(B_1\) and \(C_1\), \(A_1\) is definitely in the top 3. So we only need one final race among
\[
\{A_2, A_3, B_1, B_2, C_1\}
\]
to determine 2nd and 3rd place.

This is the 7th race.
\end{enumerate}

\begin{EditorialNote}


\textbf{Core skill tested.} This is a structured elimination problem rather than a formula dump. The interviewer is looking for disciplined bookkeeping: define the unknowns, translate each statement into a constraint, and use consistency to rule out impossible cases without losing track of the assumptions.

\textbf{How to present it in interview.} State the unknowns and the rule of the game first, then convert each sentence into a constraint you can reuse. Interviewers like hearing the consistency checks out loud because it shows that you are controlling the search space instead of guessing and backfitting.

\textbf{Common traps.}
\begin{itemize}
\item Treating each sentence in isolation instead of writing a small system of constraints.
\item Forgetting that truth/lie or ordering rules depend on who is speaking to whom.
\item Doing arithmetic too early before the structural consistency conditions are fixed.
\end{itemize}

\textbf{Related follow-ups.}
\begin{itemize}
\item Can you solve the same puzzle by contradiction or by ordering the cases differently?
\item What is the minimum set of constraints that actually determines the unique solution?
\end{itemize}


\end{EditorialNote}
\QuestionSpacer

\subsection{Classical probability and combinatorics}

\InterviewQuestionHeader{Question 3}{Options market-making and proprietary trading firm}{Quantitative Trader}{Probability / Infinite Series}{Easy}
\textbf{Problem.} Two players, A and B, alternately roll a fair six-sided die. The first player to roll a 6 wins. A goes first. What is the probability that A wins?

\textbf{Solution.}

Let \(p\) be the probability that A wins.

A wins immediately on the first roll with probability
\[
\frac{1}{6}.
\]

If both A and B fail to roll a 6 in the first round, which happens with probability
\[
\left(\frac{5}{6}\right)^2 = \frac{25}{36},
\]
then the game resets to the original state. Therefore,
\[
p = \frac{1}{6} + \frac{25}{36}p.
\]

Solving,
\[
p\left(1-\frac{25}{36}\right) = \frac{1}{6}
\quad\Rightarrow\quad
p\cdot \frac{11}{36} = \frac{1}{6}
\quad\Rightarrow\quad
p = \frac{6}{11}.
\]

Thus, the probability that A wins is \(\boxed{\frac{6}{11}}\).

\begin{EditorialNote}


\textbf{Core skill tested.} This is a structured elimination problem rather than a formula dump. The interviewer is looking for disciplined bookkeeping: define the unknowns, translate each statement into a constraint, and use consistency to rule out impossible cases without losing track of the assumptions.

\textbf{How to present it in interview.} State the unknowns and the rule of the game first, then convert each sentence into a constraint you can reuse. Interviewers like hearing the consistency checks out loud because it shows that you are controlling the search space instead of guessing and backfitting.

\textbf{Common traps.}
\begin{itemize}
\item Treating each sentence in isolation instead of writing a small system of constraints.
\item Forgetting that truth/lie or ordering rules depend on who is speaking to whom.
\item Doing arithmetic too early before the structural consistency conditions are fixed.
\end{itemize}

\textbf{Related follow-ups.}
\begin{itemize}
\item Can you solve the same puzzle by contradiction or by ordering the cases differently?
\item What is the minimum set of constraints that actually determines the unique solution?
\end{itemize}


\end{EditorialNote}
\QuestionSpacer

\InterviewQuestionHeader{Question 4}{Multi-strategy alternative investment firm}{Quantitative Trader}{Combinatorics / Probability}{Medium}
\textbf{Problem.} A fair die is rolled n times. What is the probability that the sum of all rolls equals k? i.i.d.

\textbf{Solution.}

Let \(X_1, \dots, X_n \sim \mathrm{Unif}\{1,2,3,4,5,6\}\), and let
\[
S_n = \sum_{i=1}^n X_i.
\]
Clearly, \(k \in \{n, \dots, 6n\}\); otherwise, the probability is \(0\).

The number of integer solutions to
\[
x_1 + \cdots + x_n = k, \qquad 1 \le x_i \le 6,
\]
equals the coefficient of \(z^k\) in
\[
(z + z^2 + \cdots + z^6)^n.
\]

By inclusion--exclusion,
\[
\#(k;n)
=
\sum_{j=0}^{\left\lfloor \frac{k-n}{6} \right\rfloor}
(-1)^j \binom{n}{j}\binom{k-6j-1}{n-1}.
\]

Therefore,
\[
\mathbb{P}(S_n = k)
=
\frac{1}{6^n}
\sum_{j=0}^{\left\lfloor \frac{k-n}{6} \right\rfloor}
(-1)^j \binom{n}{j}\binom{k-6j-1}{n-1},
\qquad n \le k \le 6n.
\]

\textbf{Notes.} For large \(n\), a normal approximation works:
\[
S_n \approx N\!\left(3.5n,\; \frac{35}{12}n\right),
\]
using a continuity correction.

\begin{EditorialNote}


\textbf{Core skill tested.} This is a structured elimination problem rather than a formula dump. The interviewer is looking for disciplined bookkeeping: define the unknowns, translate each statement into a constraint, and use consistency to rule out impossible cases without losing track of the assumptions.

\textbf{How to present it in interview.} State the unknowns and the rule of the game first, then convert each sentence into a constraint you can reuse. Interviewers like hearing the consistency checks out loud because it shows that you are controlling the search space instead of guessing and backfitting.

\textbf{Common traps.}
\begin{itemize}
\item Treating each sentence in isolation instead of writing a small system of constraints.
\item Forgetting that truth/lie or ordering rules depend on who is speaking to whom.
\item Doing arithmetic too early before the structural consistency conditions are fixed.
\end{itemize}

\textbf{Related follow-ups.}
\begin{itemize}
\item Can you solve the same puzzle by contradiction or by ordering the cases differently?
\item What is the minimum set of constraints that actually determines the unique solution?
\end{itemize}


\end{EditorialNote}
\QuestionSpacer

\InterviewQuestionHeader{Question 5}{Quantitative proprietary trading firm}{Quantitative Researcher}{Randomization / Fairness}{Easy}
\textbf{Problem.} Design a fair lottery system using a single biased coin with unknown bias \(p = P(H)\).

\textbf{Solution.}


Assume \(0<p<1\). Use the von Neumann extractor.

\begin{enumerate}
\item Flip the biased coin twice. If the outcome is \(HT\), output \(1\); if it is \(TH\), output \(0\); if it is \(HH\) or \(TT\), discard the pair and repeat.
\item Because
\[
P(HT)=p(1-p)=P(TH),
\]
the emitted bit is exactly fair even though \(p\) is unknown.
\item A two-flip attempt succeeds with probability \(2p(1-p)\), so the expected number of original coin flips per fair bit is
\[
\frac{2}{2p(1-p)}=\frac{1}{p(1-p)}.
\]
\end{enumerate}

To choose uniformly among \(N\) labels, generate \(k=\lceil\log_2N\rceil\) fair bits, interpret them as \(U\in\{0,\dots,2^k-1\}\), accept if \(U<N\), and otherwise reject and resample. The accepted outcome is exactly uniform on \(\{0,\dots,N-1\}\).

If \(p=0\) or \(p=1\), the coin is deterministic and no procedure using only that coin can generate a non-degenerate fair random bit.

\begin{EditorialNote}


\textbf{Core skill tested.} This is a structured elimination problem rather than a formula dump. The interviewer is looking for disciplined bookkeeping: define the unknowns, translate each statement into a constraint, and use consistency to rule out impossible cases without losing track of the assumptions.

\textbf{How to present it in interview.} State the unknowns and the rule of the game first, then convert each sentence into a constraint you can reuse. Interviewers like hearing the consistency checks out loud because it shows that you are controlling the search space instead of guessing and backfitting.

\textbf{Common traps.}
\begin{itemize}
\item Treating each sentence in isolation instead of writing a small system of constraints.
\item Forgetting that truth/lie or ordering rules depend on who is speaking to whom.
\item Doing arithmetic too early before the structural consistency conditions are fixed.
\end{itemize}

\textbf{Related follow-ups.}
\begin{itemize}
\item Can you solve the same puzzle by contradiction or by ordering the cases differently?
\item What is the minimum set of constraints that actually determines the unique solution?
\end{itemize}


\end{EditorialNote}
\QuestionSpacer

\InterviewQuestionHeader{Question 6}{Quantitative proprietary trading and market-making firm}{Quantitative Trader}{Probability}{Medium}
\textbf{Problem.} Given a fair twelve-sided die (faces 1-12), you roll repeatedly until the cumulative sum is odd. Which odd total occurs with the highest probability?

\textbf{Solution.}


Let \(p_s=P(S=s)\) for odd \(s\). On the first roll, either we roll the terminal odd value immediately, or we roll an even value \(e\in\{2,4,6,8,10,12\}\) and restart from the same problem with \(e\) already added. Hence
\[
p_s=\frac{1}{12}\mathbf 1_{\{s\in\{1,3,5,7,9,11\}\}}
+\frac1{12}\sum_{e\in\{2,4,6,8,10,12\}\atop e<s}p_{s-e}.
\]
Starting from \(p_s=0\) for \(s\le0\), this gives
\[
\begin{aligned}
p_1&=0.083333,\\
p_3&=0.090278,\\
p_5&=0.097801,\\
p_7&=0.105951,\\
p_9&=0.114780,\\
p_{11}&=0.124345.
\end{aligned}
\]
These probabilities are strictly increasing through \(11\).

For any odd \(s>11\), the direct-stop term disappears, so
\[
p_s=\frac1{12}\sum p_{s-e}\le \frac{6}{12}\max_{u<s}p_u.
\]
Thus no value above \(11\) can exceed the largest probability already attained among \(1,3,5,7,9,11\). Therefore the most likely final odd total is
\[
\boxed{11}.
\]

\begin{EditorialNote}


\textbf{Core skill tested.} This is a structured elimination problem rather than a formula dump. The interviewer is looking for disciplined bookkeeping: define the unknowns, translate each statement into a constraint, and use consistency to rule out impossible cases without losing track of the assumptions.

\textbf{How to present it in interview.} State the unknowns and the rule of the game first, then convert each sentence into a constraint you can reuse. Interviewers like hearing the consistency checks out loud because it shows that you are controlling the search space instead of guessing and backfitting.

\textbf{Common traps.}
\begin{itemize}
\item Treating each sentence in isolation instead of writing a small system of constraints.
\item Forgetting that truth/lie or ordering rules depend on who is speaking to whom.
\item Doing arithmetic too early before the structural consistency conditions are fixed.
\end{itemize}

\textbf{Related follow-ups.}
\begin{itemize}
\item Can you solve the same puzzle by contradiction or by ordering the cases differently?
\item What is the minimum set of constraints that actually determines the unique solution?
\end{itemize}


\end{EditorialNote}
\QuestionSpacer

\InterviewQuestionHeader{Question 7}{Quantitative proprietary trading and market-making firm}{Quantitative Trader}{Probability}{Medium}
\textbf{Problem.} A stick is broken at one random point. Then one of the pieces is chosen at random and broken at a random point. What is the probability that the three pieces can form a triangle?

\textbf{Solution.}

Let the first break be at
\[
X \sim \mathrm{Unif}(0,1),
\]
creating pieces of lengths \(X\) and \(1-X\). Choose one of these two pieces with probability \(1/2\); call its length
\[
L \in \{X,1-X\}.
\]
Now break that chosen piece at
\[
U \sim \mathrm{Unif}(0,1),
\]
producing lengths
\[
UL,\qquad (1-U)L,\qquad 1-L.
\]

Since the total length is \(1\), the three pieces form a triangle if and only if no piece has length at least \(1/2\), i.e.
\[
\max\{UL,(1-U)L,1-L\}<\frac12.
\]

\textbf{Case 1:} choose the \(X\)-piece (probability \(1/2\)).

We need
\[
1-X<\frac12 \qquad\Longrightarrow\qquad X>\frac12,
\]
and also
\[
UX<\frac12,\qquad (1-U)X<\frac12.
\]
These give
\[
U<\frac{1}{2X},
\qquad
U>1-\frac{1}{2X}.
\]
So for fixed \(X>\tfrac12\), the allowed interval for \(U\) is
\[
1-\frac{1}{2X}<U<\frac{1}{2X},
\]
whose length is
\[
\frac{1}{X}-1.
\]
Thus the contribution from this case is
\[
\frac12 \int_{1/2}^1 \left(\frac{1}{x}-1\right)\,dx
=
\frac12\bigl(\ln 2-\tfrac12\bigr).
\]

\textbf{Case 2:} choose the \((1-X)\)-piece (probability \(1/2\)).

By symmetry, this requires
\[
X<\frac12,
\]
and the allowed interval for \(U\) has length
\[
\frac{1}{1-X}-1.
\]
Hence the contribution is
\[
\frac12 \int_0^{1/2} \left(\frac{1}{1-x}-1\right)\,dx
=
\frac12\bigl(\ln 2-\tfrac12\bigr).
\]

Adding the two cases,
\[
\mathbb{P}(\text{triangle})
=
\ln 2-\frac12
\approx 0.1931.
\]

\textbf{Note.} This differs from the classic problem of making two independent uniform cuts on the whole stick, which gives probability \(1/4\).

\begin{EditorialNote}


\textbf{Core skill tested.} This is a structured elimination problem rather than a formula dump. The interviewer is looking for disciplined bookkeeping: define the unknowns, translate each statement into a constraint, and use consistency to rule out impossible cases without losing track of the assumptions.

\textbf{How to present it in interview.} State the unknowns and the rule of the game first, then convert each sentence into a constraint you can reuse. Interviewers like hearing the consistency checks out loud because it shows that you are controlling the search space instead of guessing and backfitting.

\textbf{Common traps.}
\begin{itemize}
\item Treating each sentence in isolation instead of writing a small system of constraints.
\item Forgetting that truth/lie or ordering rules depend on who is speaking to whom.
\item Doing arithmetic too early before the structural consistency conditions are fixed.
\end{itemize}

\textbf{Related follow-ups.}
\begin{itemize}
\item Can you solve the same puzzle by contradiction or by ordering the cases differently?
\item What is the minimum set of constraints that actually determines the unique solution?
\end{itemize}


\end{EditorialNote}
\QuestionSpacer

\subsection{Conditional probability, Bayes, and expectation}

\InterviewQuestionHeader{Question 8}{Quantitative proprietary trading and market-making firm}{Quantitative Trader}{Probability / Bayes}{Easy}
\textbf{Problem.} You have 100 coins: 99 are fair and 1 is double-headed. You pick one uniformly at random and flip it 10 times, observing 10 heads. What is the probability that you picked the double-headed coin?

\textbf{Solution.}

Let \(D\) be the event ``picked the double-headed coin,'' and let \(H\) be the event ``10 heads in 10 flips.''

The prior probabilities are
\[
\mathbb{P}(D)=\frac{1}{100},
\qquad
\mathbb{P}(D^c)=\frac{99}{100}.
\]
The likelihoods are
\[
\mathbb{P}(H\mid D)=1,
\qquad
\mathbb{P}(H\mid D^c)=\left(\frac12\right)^{10}=\frac{1}{1024}.
\]

By Bayes' rule,
\[
\mathbb{P}(D\mid H)
=
\frac{\mathbb{P}(H\mid D)\mathbb{P}(D)}
{\mathbb{P}(H\mid D)\mathbb{P}(D)+\mathbb{P}(H\mid D^c)\mathbb{P}(D^c)}.
\]
Substituting,
\[
\mathbb{P}(D\mid H)
=
\frac{1\cdot \frac{1}{100}}
{1\cdot \frac{1}{100}+\frac{1}{1024}\cdot \frac{99}{100}}
=
\frac{1}{1+\frac{99}{1024}}
=
\frac{1024}{1123}
\approx 0.912.
\]

Therefore, the probability that you picked the double-headed coin is
\[
\boxed{\frac{1024}{1123}\approx 0.912}.
\]

\begin{EditorialNote}


\textbf{Core skill tested.} This is a pattern-recognition problem. A strong candidate identifies the smallest correct framework first, writes the governing event or quantity cleanly, and only then performs the arithmetic or algebra needed to finish the argument.

\textbf{How to present it in interview.} In interview conditions, announce the trick early. Say whether you are using symmetry, conditioning, a complement, an invariant, or a replication argument. That helps the interviewer follow your reasoning and reduces the chance that arithmetic clutter hides the main idea.

\textbf{Common traps.}
\begin{itemize}
\item Rushing into arithmetic before defining the relevant event or random variable.
\item Missing a tie rule, support restriction, or conditioning case that changes the answer.
\item Failing to sanity-check the final value against an obvious bound or symmetry argument.
\end{itemize}

\textbf{Related follow-ups.}
\begin{itemize}
\item What is the shortest alternative derivation: symmetry, complement, recursion, or indicator variables?
\item How would the answer change under a biased draw, more players, or a different support?
\end{itemize}


\end{EditorialNote}
\QuestionSpacer

\section{Statistics, Machine Learning, and Econometrics}
\ChapterIntro{Use these questions to practice complete solutions, concise explanations, and likely follow-up questions.}

\subsection{Applied ML, feature engineering, and data quality}

\InterviewQuestionHeader{Question 9}{Global investment bank and financial services firm}{Risk / Quant}{Data Quality}{Easy}
\textbf{Problem.} Give a step-by-step checklist for cleaning and validating financial data.

\textbf{Solution.}

\begin{enumerate}
\item \textbf{Schema and keys:} check uniqueness, referential integrity, and consistency of units and currencies.
\item \textbf{Time alignment:} verify calendars, time zones, holiday rolls, and point-in-time reference data.
\item \textbf{Missing values and outliers:} define treatment rules, apply winsorization where appropriate, and adjust for corporate actions such as splits and dividends.
\item \textbf{Reconciliation:} compare source data against the golden copy, verify roll-forwards, and reconcile totals versus subtotals.
\item \textbf{Testing and governance:} run deterministic unit tests, statistical sanity checks, and maintain an audit trail with versioning.
\end{enumerate}

\begin{EditorialNote}


\textbf{Core skill tested.} This is a pattern-recognition problem. A strong candidate identifies the smallest correct framework first, writes the governing event or quantity cleanly, and only then performs the arithmetic or algebra needed to finish the argument.

\textbf{How to present it in interview.} In interview conditions, announce the trick early. Say whether you are using symmetry, conditioning, a complement, an invariant, or a replication argument. That helps the interviewer follow your reasoning and reduces the chance that arithmetic clutter hides the main idea.

\textbf{Common traps.}
\begin{itemize}
\item Rushing into arithmetic before defining the relevant event or random variable.
\item Missing a tie rule, support restriction, or conditioning case that changes the answer.
\item Failing to sanity-check the final value against an obvious bound or symmetry argument.
\end{itemize}

\textbf{Related follow-ups.}
\begin{itemize}
\item What is the shortest alternative derivation: symmetry, complement, recursion, or indicator variables?
\item How would the answer change under a biased draw, more players, or a different support?
\end{itemize}


\end{EditorialNote}
\QuestionSpacer

\InterviewQuestionHeader{Question 10}{Quantitative investment management firm}{Quantitative Researcher}{Data Preprocessing / Missing Data}{Medium}
\textbf{Problem.} You have a dataset with missing values. How would you handle this for a machine-learning model?

\textbf{Solution.}


Start by diagnosing the missingness mechanism (MCAR, MAR, or plausibly MNAR) and whether missingness itself contains signal. Then build the treatment inside a leakage-safe validation pipeline.

\begin{enumerate}
\item \textbf{Profile missingness.} Measure rates by feature/time/entity, examine missingness indicators, and test whether missingness is associated with observed covariates or the target.
\item \textbf{Choose a method.} Simple median/mode imputation is a strong baseline; MICE, k-NN, EM, or matrix-completion methods may help when justified; many tree models can handle missing values directly.
\item \textbf{Respect time.} For forecasting, only use information that would have existed at the prediction timestamp. Forward fill can be valid with an age cap when the last observation is genuinely known; backward fill generally leaks future information and should not cross a forecasting boundary.
\item \textbf{Fit inside each fold.} Estimate imputation parameters using training data only. Never compute medians, nearest neighbors, or learned imputers using validation/test observations.
\item \textbf{Validate uncertainty and production behavior.} Compare with/without missingness indicators, monitor missing-rate drift, define all-missing fallbacks, and for multiple imputation propagate between-imputation uncertainty when inferential validity matters.
\end{enumerate}
Under MNAR, no imputation algorithm automatically removes bias; sensitivity analysis or an explicit missingness model may be required.

\begin{EditorialNote}


\textbf{Core skill tested.} This item tests whether you can move cleanly from definition to estimation to validation. A polished response states the estimator or framework, clarifies its assumptions, and then points out the most likely failure mode such as leakage, non-stationarity, or unstable inference.

\textbf{How to present it in interview.} A concise interview answer should always separate modelling from evaluation. State the model or diagnostic, then immediately describe the validation design and the reason you trust or distrust the result under realistic data conditions.

\textbf{Common traps.}
\begin{itemize}
\item Stating an estimator but not the assumptions needed for interpretation.
\item Forgetting to separate train/validation/test logic or time ordering.
\item Choosing a metric that does not match the business or modelling objective.
\end{itemize}

\textbf{Related follow-ups.}
\begin{itemize}
\item What breaks first if the data are non-stationary, heavy-tailed, or missing not at random?
\item How would you redesign the validation scheme if the labels overlap or the data are strictly time ordered?
\end{itemize}


\end{EditorialNote}
\QuestionSpacer

\subsection{Linear models, estimation, and linear algebra}

\InterviewQuestionHeader{Question 11}{Cross-industry quantitative-finance employers}{Quant Research}{OLS / Linear Algebra}{Easy}
\textbf{Problem.} What is the closed-form ordinary least squares (OLS) estimator for \ensuremath{\beta} in y = X\ensuremath{\beta} + \ensuremath{\varepsilon}?

\textbf{Solution.}

If
\[
X \in \mathbb{R}^{n\times d}
\]
has full column rank (\(\operatorname{rank}(X)=d\le n\)), then the normal equations are
\[
X^\top X\,\hat{\beta}=X^\top y.
\]
Since \(X^\top X\) is invertible in this case,
\[
\hat{\beta}=(X^\top X)^{-1}X^\top y.
\]

If \(X^\top X\) is singular, use the Moore--Penrose pseudoinverse:
\[
\hat{\beta}=X^+ y.
\]

\begin{EditorialNote}


\textbf{Core skill tested.} This is a structured elimination problem rather than a formula dump. The interviewer is looking for disciplined bookkeeping: define the unknowns, translate each statement into a constraint, and use consistency to rule out impossible cases without losing track of the assumptions.

\textbf{How to present it in interview.} State the unknowns and the rule of the game first, then convert each sentence into a constraint you can reuse. Interviewers like hearing the consistency checks out loud because it shows that you are controlling the search space instead of guessing and backfitting.

\textbf{Common traps.}
\begin{itemize}
\item Treating each sentence in isolation instead of writing a small system of constraints.
\item Forgetting that truth/lie or ordering rules depend on who is speaking to whom.
\item Doing arithmetic too early before the structural consistency conditions are fixed.
\end{itemize}

\textbf{Related follow-ups.}
\begin{itemize}
\item Can you solve the same puzzle by contradiction or by ordering the cases differently?
\item What is the minimum set of constraints that actually determines the unique solution?
\end{itemize}


\end{EditorialNote}
\QuestionSpacer

\InterviewQuestionHeader{Question 12}{Cross-industry quantitative-finance employers}{Quant Research}{Regression / Multicollinearity}{Medium}
\textbf{Problem.} Explain multicollinearity, how to detect it, and how ridge regression remedies it.

\textbf{Solution.}

\textbf{Definition.} Multicollinearity occurs when predictors are exactly or nearly linearly dependent. This inflates the variance of \(\hat{\beta}\) and makes
\[
X^\top X
\]
ill-conditioned or singular.

\textbf{Detection.}
\begin{itemize}
\item inspect pairwise correlation heatmaps;
\item regress each predictor \(X_i\) on the others and compute the variance inflation factor
\[
\mathrm{VIF}_i=\frac{1}{1-R_i^2};
\]
large VIF values (for example, greater than \(5\) or \(10\)) indicate potential multicollinearity;
\item check the condition number or the eigenvalues of \(X^\top X\).
\end{itemize}

\textbf{Ridge remedy.} Ridge regression solves
\[
\hat{\beta}_\lambda
=
\arg\min_{\beta}\ \|y-X\beta\|_2^2+\lambda \|\beta\|_2^2.
\]
The first-order condition gives
\[
(X^\top X+\lambda I)\hat{\beta}_\lambda=X^\top y,
\]
so
\[
\hat{\beta}_\lambda=(X^\top X+\lambda I)^{-1}X^\top y.
\]

Adding \(\lambda I\) makes the system strictly positive definite for \(\lambda>0\), stabilizes inversion, and shrinks coefficient estimates toward zero, thereby reducing variance.

\begin{EditorialNote}


\textbf{Core skill tested.} This item tests whether you can move cleanly from definition to estimation to validation. A polished response states the estimator or framework, clarifies its assumptions, and then points out the most likely failure mode such as leakage, non-stationarity, or unstable inference.

\textbf{How to present it in interview.} A concise interview answer should always separate modelling from evaluation. State the model or diagnostic, then immediately describe the validation design and the reason you trust or distrust the result under realistic data conditions.

\textbf{Common traps.}
\begin{itemize}
\item Stating an estimator but not the assumptions needed for interpretation.
\item Forgetting to separate train/validation/test logic or time ordering.
\item Choosing a metric that does not match the business or modelling objective.
\end{itemize}

\textbf{Related follow-ups.}
\begin{itemize}
\item What breaks first if the data are non-stationary, heavy-tailed, or missing not at random?
\item How would you redesign the validation scheme if the labels overlap or the data are strictly time ordered?
\end{itemize}


\end{EditorialNote}
\QuestionSpacer

\subsection{Time series, stationarity, and diagnostics}

\InterviewQuestionHeader{Question 13}{Quantitative investment management firm}{Quantitative Researcher}{Time Series / Unit-Root Testing}{Medium}
\textbf{Problem.} You observe a time series that appears to have a unit root. How would you test for stationarity and what are the implications for modeling?

\textbf{Solution.}

\textbf{Step 1 --- Visual and preliminary checks.} Inspect the ACF/PACF and rolling summary statistics. If the variance drifts or the ACF decays slowly, a unit root is plausible.

\textbf{Step 2 --- Augmented Dickey--Fuller (ADF) test.} Estimate
\[
\Delta y_t = \alpha + \beta t + \gamma y_{t-1} + \sum_{i=1}^{p}\phi_i \Delta y_{t-i} + \varepsilon_t,
\]
and test
\[
H_0:\gamma=0 \qquad \text{(unit root, \(y_t\) is \(I(1)\))},
\]
versus
\[
H_1:\gamma<0 \qquad \text{(stationary)}.
\]
Choose the lag \(p\) by AIC or BIC so that \(\varepsilon_t\) is approximately white noise, and include the trend term \(\beta t\) only if the series level exhibits a deterministic trend. The ADF statistic uses non-standard Dickey--Fuller critical values.

\textbf{Step 3 --- Complementary checks.}

\begin{itemize}
\item \textbf{KPSS test:} the null is (trend-)stationarity. Rejection together with ADF non-rejection strengthens evidence of a unit root.
\item \textbf{Phillips--Perron (PP) test:} robust to more general serial correlation and heteroskedasticity.
\end{itemize}

\textbf{Implications for modeling.}

\begin{itemize}
\item If a unit root is confirmed, difference the series until it is stationary and model it with ARIMA\((p,d,q)\), where \(d\) is the differencing order.
\item If multiple series are \(I(1)\) but cointegrated, use a vector error-correction model (VECM):
\[
\Delta y_t = \Pi y_{t-1} + \sum_{i=1}^{k-1}\Gamma_i \Delta y_{t-i} + u_t,
\qquad
\Pi=\alpha\beta^\top,
\]
or a single-equation ECM of the form
\[
\Delta y_t = \lambda (y_{t-1}-\beta^\top x_{t-1}) + \cdots.
\]
\item Forecasting and inference should be done on stationary transforms; level forecasts can then be re-integrated.
\end{itemize}

\begin{EditorialNote}


\textbf{Core skill tested.} Geometric probability questions are usually won by choosing the right picture before doing any algebra. The interviewer wants to see whether you can translate the verbal statement into a region, a symmetry argument, or a measure-preserving transform and then compute only what is necessary.

\textbf{How to present it in interview.} In interview conditions, announce the trick early. Say whether you are using symmetry, conditioning, a complement, an invariant, or a replication argument. That helps the interviewer follow your reasoning and reduces the chance that arithmetic clutter hides the main idea.

\textbf{Common traps.}
\begin{itemize}
\item Rushing into arithmetic before defining the relevant event or random variable.
\item Missing a tie rule, support restriction, or conditioning case that changes the answer.
\item Failing to sanity-check the final value against an obvious bound or symmetry argument.
\end{itemize}

\textbf{Related follow-ups.}
\begin{itemize}
\item What is the shortest alternative derivation: symmetry, complement, recursion, or indicator variables?
\item How would the answer change under a biased draw, more players, or a different support?
\end{itemize}


\end{EditorialNote}
\QuestionSpacer

\section{Derivatives, Stochastic Calculus, and Option Theory}
\ChapterIntro{Use these questions to practice complete solutions, concise explanations, and likely follow-up questions.}

\subsection{GBM, Ito calculus, martingales, and Monte Carlo}

\InterviewQuestionHeader{Question 14}{Multi-strategy alternative investment firm}{Quantitative Trader}{Stochastic Calculus}{Medium}
\textbf{Problem.} Derive the closed-form solution of geometric Brownian motion (GBM) and briefly explain its main applications in finance.

\textbf{Solution.}

Let \(S_t\) follow
\[
\frac{dS_t}{S_t} = \mu\,dt + \sigma\,dW_t,
\qquad
S_0>0,
\]
with constant drift \(\mu\), volatility \(\sigma>0\), and \((W_t)_{t\ge 0}\) a standard Brownian motion.

Apply Itô's lemma to \(f(S_t)=\ln S_t\). Then
\[
d\ln S_t
=
\frac{1}{S_t}\,dS_t-\frac{1}{2S_t^2}(dS_t)^2
=
\left(\mu-\frac12\sigma^2\right)dt+\sigma\,dW_t.
\]

Integrating from \(0\) to \(t\),
\[
\ln S_t
=
\ln S_0+\left(\mu-\frac12\sigma^2\right)t+\sigma W_t.
\]
Therefore,
\[
S_t
=
S_0\exp\!\left(\left(\mu-\frac12\sigma^2\right)t+\sigma W_t\right).
\]

Hence
\[
\ln S_t \sim N\!\left(\ln S_0+\left(\mu-\frac12\sigma^2\right)t,\ \sigma^2 t\right),
\]
so \(S_t\) is lognormal.

Its moments are
\[
\mathbb{E}[S_t]=S_0 e^{\mu t},
\]
and
\[
\mathrm{Var}(S_t)=S_0^2 e^{2\mu t}\left(e^{\sigma^2 t}-1\right).
\]

\textbf{Applications in finance.}

\begin{itemize}
\item \textbf{Derivative pricing (Black--Scholes):} under the risk-neutral measure, replace \(\mu\) by \(r-q\), where \(r\) is the risk-free rate and \(q\) is the dividend yield. The GBM solution gives the lognormal terminal distribution used in Black--Scholes and Monte Carlo pricing.
\item \textbf{Risk management:} simulate P\&L scenarios for equities, FX, and commodities; compute risk metrics such as VaR and expected shortfall under lognormal dynamics.
\item \textbf{Volatility and elasticity intuition:} log-returns are normal with variance \(\sigma^2 t\), and relative changes scale with price, which avoids negative prices.
\item \textbf{Portfolio problems:} continuous-time growth-optimal/Kelly analysis and Merton portfolio choice with GBM assets.
\end{itemize}

\begin{EditorialNote}


\textbf{Core skill tested.} The core skill here is to connect a formula or payoff identity to the economic reason it matters. Good answers do not stop at the equation; they also explain which no-arbitrage argument, SDE manipulation, or hedging interpretation makes the result useful in practice.

\textbf{How to present it in interview.} The clean interview version is usually a 30-to-60 second explanation: define the process or parity, write the one essential equation, and state what it implies for pricing, hedging, or risk. Avoid drowning the answer in notation unless the interviewer explicitly asks for the full derivation.

\textbf{Common traps.}
\begin{itemize}
\item Mixing up payoff, profit, and present value, or confusing physical and risk-neutral interpretations.
\item Using the right formula but never saying what assumption makes it valid.
\item Losing points on sign conventions, discounting, or boundary conditions.
\end{itemize}

\textbf{Related follow-ups.}
\begin{itemize}
\item How would the answer change under dividends, stochastic volatility, or a different compounding convention?
\item Can you derive the same result from replication, PDE language, and intuition about convexity or hedging?
\end{itemize}


\end{EditorialNote}
\QuestionSpacer

\InterviewQuestionHeader{Question 15}{Multi-strategy alternative investment firm}{Quantitative Trader}{Stochastic Calculus / Derivatives}{Medium}
\textbf{Problem.} What is Itô's lemma and how is it used in derivative pricing?

\textbf{Solution.}

\textbf{Itô's lemma (one-dimensional diffusion).} Let \(X_t\) satisfy
\[
dX_t = a(t,X_t)\,dt + b(t,X_t)\,dW_t,
\]
and let \(f(t,x)\in C^{1,2}\). Then
\[
df(t,X_t)
=
f_t\,dt + f_x\,dX_t + \frac12 f_{xx}(dX_t)^2.
\]
Using
\[
(dW_t)^2=dt,\qquad dt\,dW_t=0,\qquad dt^2=0,
\]
this becomes
\[
df(t,X_t)
=
\left(f_t + a f_x + \frac12 b^2 f_{xx}\right)dt + b f_x\,dW_t.
\]

\textbf{Use in risk-neutral derivative pricing.}

Suppose the underlying follows GBM under the risk-neutral measure \(\mathbb Q\):
\[
dS_t = rS_t\,dt + \sigma S_t\,dW_t^{\mathbb Q}.
\]
For a derivative price \(V(t,S)=f(t,S)\), apply Itô's lemma:
\[
dV
=
V_t\,dt + V_S\,dS + \frac12 V_{SS}(dS)^2.
\]
Substituting \(dS\),
\[
dV
=
\left(V_t + rS V_S + \frac12 \sigma^2 S^2 V_{SS}\right)dt
+
\sigma S V_S\,dW_t^{\mathbb Q}.
\]

Now form a self-financing delta-hedged portfolio
\[
\Pi = V - \Delta S,
\qquad
\Delta = V_S.
\]
The \(dW_t^{\mathbb Q}\) term cancels, so \(\Pi\) is locally riskless and must earn the risk-free rate:
\[
d\Pi = r\Pi\,dt.
\]
Hence
\[
V_t + rS V_S + \frac12 \sigma^2 S^2 V_{SS} - rV = 0,
\]
with terminal condition
\[
V(T,S)=\text{payoff}(S).
\]

This is the Black--Scholes PDE. Solving it, or equivalently taking the risk-neutral expectation, gives derivative prices. For example, with payoff
\[
V(T,S)=(S-K)^+,
\]
the solution is the Black--Scholes call formula
\[
C = S N(d_1) - K e^{-r(T-t)} N(d_2).
\]

\textbf{Remarks.} Itô's lemma is also used to:
\begin{itemize}
\item derive the lognormality of GBM via \(f=\ln S\);
\item obtain Greeks such as \(V_S\) and \(V_{SS}\);
\item extend pricing arguments to multi-factor models via the multidimensional Itô formula with cross terms.
\end{itemize}

\begin{EditorialNote}


\textbf{Core skill tested.} The core skill here is to connect a formula or payoff identity to the economic reason it matters. Good answers do not stop at the equation; they also explain which no-arbitrage argument, SDE manipulation, or hedging interpretation makes the result useful in practice.

\textbf{How to present it in interview.} The clean interview version is usually a 30-to-60 second explanation: define the process or parity, write the one essential equation, and state what it implies for pricing, hedging, or risk. Avoid drowning the answer in notation unless the interviewer explicitly asks for the full derivation.

\textbf{Common traps.}
\begin{itemize}
\item Mixing up payoff, profit, and present value, or confusing physical and risk-neutral interpretations.
\item Using the right formula but never saying what assumption makes it valid.
\item Losing points on sign conventions, discounting, or boundary conditions.
\end{itemize}

\textbf{Related follow-ups.}
\begin{itemize}
\item How would the answer change under dividends, stochastic volatility, or a different compounding convention?
\item Can you derive the same result from replication, PDE language, and intuition about convexity or hedging?
\end{itemize}


\end{EditorialNote}
\QuestionSpacer

\InterviewQuestionHeader{Question 16}{Global investment bank and financial services firm}{Risk / Quant}{Monte Carlo}{Medium}
\textbf{Problem.} Describe the Monte Carlo workflow for market risk.

\textbf{Solution.}

\begin{enumerate}
\item Specify marginals and dependence, for example with a Gaussian copula or dynamic volatility model.
\item Generate \(M\) scenarios, or full paths if the problem is path-dependent.
\item Reprice positions under each scenario, using either full revaluation or a risk-factor approximation.
\item Aggregate losses and compute metrics such as VaR, expected shortfall, PFE, and CVA.
\item Improve efficiency with antithetic variates, control variates, stratified sampling, or quasi-random sequences, and perform convergence checks.
\end{enumerate}

\begin{EditorialNote}


\textbf{Core skill tested.} This is a risk-and-finance framing problem. The best answer defines the object precisely, states the minimal formula or workflow, and then immediately mentions the practical assumption set, because interviewers care about whether you know where textbook simplifications break in a desk setting.

\textbf{How to present it in interview.} Answer these by sequencing definition -> construction -> caveat. For example, define the measure or sensitivity, say how you would compute it on real data or positions, and then mention one limitation such as liquidity, model error, or convention dependence.

\textbf{Common traps.}
\begin{itemize}
\item Quoting a definition without specifying horizon, confidence level, or market convention.
\item Ignoring model risk, liquidity, or scenario design when moving from textbook to practice.
\item Giving a sensitivity or valuation formula without saying what inputs must be bootstrapped or estimated.
\end{itemize}

\textbf{Related follow-ups.}
\begin{itemize}
\item What part of the workflow is convention-dependent, and what part is model-dependent?
\item How would you explain the same concept to a portfolio manager instead of a quant interviewer?
\end{itemize}


\end{EditorialNote}
\QuestionSpacer

\subsection{Options, no-arbitrage, Black-Scholes, Greeks, and hedging}

\InterviewQuestionHeader{Question 17}{Global investment bank and financial services firm}{Risk / Quant}{Black–Scholes}{Easy}
\textbf{Problem.} Summarize the Black-Scholes framework, key assumptions, and pricing steps.

\textbf{Solution.}


Under Black--Scholes--Merton with continuous dividend yield \(q\),
\[
\frac{dS_t}{S_t}=(r-q)\,dt+\sigma\,dW_t^{\mathbb Q}
\]
under the risk-neutral measure. The model assumes no arbitrage, frictionless continuous trading in the idealization, constant \(r,q,\sigma\), and lognormal diffusion dynamics.

For a European call,
\[
C=S_0e^{-qT}\Phi(d_1)-Ke^{-rT}\Phi(d_2),
\]
where
\[
d_1=\frac{\ln(S_0/K)+(r-q+\tfrac12\sigma^2)T}{\sigma\sqrt T},
\qquad
d_2=d_1-\sigma\sqrt T.
\]
The put price follows from
\[
P=Ke^{-rT}\Phi(-d_2)-S_0e^{-qT}\Phi(-d_1)
\]
or from put--call parity
\[
C-P=S_0e^{-qT}-Ke^{-rT}.
\]
For a non-dividend-paying stock set \(q=0\).

Implementation steps: validate inputs and units, compute \(d_1,d_2\), evaluate the normal CDF, price the option, and sanity-check intrinsic-value bounds, limiting cases, and put--call parity. The major practical caveat is that market implied volatility varies by strike and maturity, so constant \(\sigma\) is an idealization rather than an empirical surface model.

\begin{EditorialNote}


\textbf{Core skill tested.} The core skill here is to connect a formula or payoff identity to the economic reason it matters. Good answers do not stop at the equation; they also explain which no-arbitrage argument, SDE manipulation, or hedging interpretation makes the result useful in practice.

\textbf{How to present it in interview.} The clean interview version is usually a 30-to-60 second explanation: define the process or parity, write the one essential equation, and state what it implies for pricing, hedging, or risk. Avoid drowning the answer in notation unless the interviewer explicitly asks for the full derivation.

\textbf{Common traps.}
\begin{itemize}
\item Mixing up payoff, profit, and present value, or confusing physical and risk-neutral interpretations.
\item Using the right formula but never saying what assumption makes it valid.
\item Losing points on sign conventions, discounting, or boundary conditions.
\end{itemize}

\textbf{Related follow-ups.}
\begin{itemize}
\item How would the answer change under dividends, stochastic volatility, or a different compounding convention?
\item Can you derive the same result from replication, PDE language, and intuition about convexity or hedging?
\end{itemize}


\end{EditorialNote}
\QuestionSpacer

\InterviewQuestionHeader{Question 18}{Global investment bank and financial services firm}{Risk / Quant}{Greeks}{Easy}
\textbf{Problem.} Define the main Greeks and how you would use them in hedging.

\textbf{Solution.}


Let \(\tau\) denote \emph{time to expiry}. For a Black--Scholes--Merton call with dividend yield \(q\), the main Greeks include
\[
\Delta=\frac{\partial C}{\partial S}=e^{-q\tau}\Phi(d_1),
\]
\[
\Gamma=\frac{\partial^2 C}{\partial S^2}
=\frac{e^{-q\tau}\phi(d_1)}{S\sigma\sqrt\tau},
\]
\[
\mathrm{Vega}=\frac{\partial C}{\partial\sigma}
=Se^{-q\tau}\phi(d_1)\sqrt\tau,
\]
and
\[
\rho=\frac{\partial C}{\partial r}=K\tau e^{-r\tau}\Phi(d_2).
\]

Be explicit about the theta convention. Market \emph{theta} is usually the derivative with respect to calendar time \(t\) while expiry date is fixed:
\[
\Theta=\frac{\partial C}{\partial t}=-\frac{\partial C}{\partial\tau}.
\]
So calling \(\partial C/\partial\tau\) ``time decay'' without the minus sign reverses the usual convention.

In hedging, use the underlying to neutralize delta, options to control gamma and vega, and rate instruments when rho matters. Rebalance according to time/trigger rules and measure residual risk, transaction costs, discrete-hedging error, skew/term-structure exposure, and model risk.

\begin{EditorialNote}


\textbf{Core skill tested.} The core skill here is to connect a formula or payoff identity to the economic reason it matters. Good answers do not stop at the equation; they also explain which no-arbitrage argument, SDE manipulation, or hedging interpretation makes the result useful in practice.

\textbf{How to present it in interview.} The clean interview version is usually a 30-to-60 second explanation: define the process or parity, write the one essential equation, and state what it implies for pricing, hedging, or risk. Avoid drowning the answer in notation unless the interviewer explicitly asks for the full derivation.

\textbf{Common traps.}
\begin{itemize}
\item Mixing up payoff, profit, and present value, or confusing physical and risk-neutral interpretations.
\item Using the right formula but never saying what assumption makes it valid.
\item Losing points on sign conventions, discounting, or boundary conditions.
\end{itemize}

\textbf{Related follow-ups.}
\begin{itemize}
\item How would the answer change under dividends, stochastic volatility, or a different compounding convention?
\item Can you derive the same result from replication, PDE language, and intuition about convexity or hedging?
\end{itemize}


\end{EditorialNote}
\QuestionSpacer

\section{Risk, Fixed Income, and Portfolio Management}
\ChapterIntro{Use these questions to practice complete solutions, concise explanations, and likely follow-up questions.}

\subsection{Risk measures, stress testing, and volatility}

\InterviewQuestionHeader{Question 19}{Global investment bank and financial services firm}{Risk / Quant}{Value at Risk}{Easy}
\textbf{Problem.} Explain Value at Risk (VaR) and its limitations. Provide a step-by-step way to compute it.

\textbf{Solution.}


\begin{enumerate}
\item \textbf{Definition.} If \(L\) denotes portfolio loss, VaR at confidence level \(\alpha\) is the \(\alpha\)-quantile of the loss distribution:
\[
\mathrm{VaR}_\alpha=\inf\{\ell:P(L\le\ell)\ge\alpha\}.
\]
For a continuous distribution this is equivalent to \(P(L>\mathrm{VaR}_\alpha)=1-\alpha\).

\item \textbf{Historical simulation.} Revalue the current portfolio under historical shocks (or use horizon-matched historical P\&L), convert results to losses, sort losses from smallest to largest, and take the empirical \(\alpha\)-quantile. Thus 99\% VaR is the 99th percentile of loss, not the 1st percentile.

\item \textbf{Delta-normal VaR.} If
\[
\Delta P\approx \Delta^\top\Delta X,
\qquad \Delta X\sim N(0,\Sigma),
\]
then the loss standard deviation is \(\sqrt{\Delta^\top\Sigma\Delta}\), and for zero mean
\[
\mathrm{VaR}_\alpha=\Phi^{-1}(\alpha)\sqrt{\Delta^\top\Sigma\Delta}.
\]
Add the mean term if it is material and use delta-gamma or full revaluation for nonlinear portfolios.

\item \textbf{Monte Carlo VaR.} Simulate joint risk-factor scenarios, reprice the portfolio, convert P\&L to loss, and again take the empirical \(\alpha\)-quantile of loss.
\end{enumerate}

Limitations include model and parameter risk, liquidity assumptions, procyclicality, and the fact that VaR says nothing about the average severity beyond the quantile. Expected Shortfall addresses the last issue and is coherent under standard conditions, whereas VaR need not be subadditive.

\begin{EditorialNote}


\textbf{Core skill tested.} This is a risk-and-finance framing problem. The best answer defines the object precisely, states the minimal formula or workflow, and then immediately mentions the practical assumption set, because interviewers care about whether you know where textbook simplifications break in a desk setting.

\textbf{How to present it in interview.} Answer these by sequencing definition -> construction -> caveat. For example, define the measure or sensitivity, say how you would compute it on real data or positions, and then mention one limitation such as liquidity, model error, or convention dependence.

\textbf{Common traps.}
\begin{itemize}
\item Quoting a definition without specifying horizon, confidence level, or market convention.
\item Ignoring model risk, liquidity, or scenario design when moving from textbook to practice.
\item Giving a sensitivity or valuation formula without saying what inputs must be bootstrapped or estimated.
\end{itemize}

\textbf{Related follow-ups.}
\begin{itemize}
\item What part of the workflow is convention-dependent, and what part is model-dependent?
\item How would you explain the same concept to a portfolio manager instead of a quant interviewer?
\end{itemize}


\end{EditorialNote}
\QuestionSpacer

\InterviewQuestionHeader{Question 20}{Global investment bank and financial services firm}{Risk / Quant}{Stress Testing}{Medium}
\textbf{Problem.} How do you design stress scenarios?

\textbf{Solution.}

\begin{enumerate}
\item \textbf{Scenario type.} Use historical replay (for example, the GFC), hypothetical shocks, and reverse stress tests to identify breaking points.

\item \textbf{Scenario construction.} Apply joint shocks that respect cross-asset relationships, basis effects, and liquidity horizons.

\item \textbf{Calibration.} Anchor severity to the firm's risk appetite, historical precedent, and regulatory requirements.

\item \textbf{Outputs.} Report P\&L by factor and desk, breached limits, management actions, and governance sign-off.
\end{enumerate}

\begin{EditorialNote}


\textbf{Core skill tested.} This is a risk-and-finance framing problem. The best answer defines the object precisely, states the minimal formula or workflow, and then immediately mentions the practical assumption set, because interviewers care about whether you know where textbook simplifications break in a desk setting.

\textbf{How to present it in interview.} Answer these by sequencing definition -> construction -> caveat. For example, define the measure or sensitivity, say how you would compute it on real data or positions, and then mention one limitation such as liquidity, model error, or convention dependence.

\textbf{Common traps.}
\begin{itemize}
\item Quoting a definition without specifying horizon, confidence level, or market convention.
\item Ignoring model risk, liquidity, or scenario design when moving from textbook to practice.
\item Giving a sensitivity or valuation formula without saying what inputs must be bootstrapped or estimated.
\end{itemize}

\textbf{Related follow-ups.}
\begin{itemize}
\item What part of the workflow is convention-dependent, and what part is model-dependent?
\item How would you explain the same concept to a portfolio manager instead of a quant interviewer?
\end{itemize}


\end{EditorialNote}
\QuestionSpacer

\subsection{Credit and counterparty risk}

\InterviewQuestionHeader{Question 21}{Global investment bank and financial services firm}{Risk / Quant}{Credit Risk Modelling}{Medium}
\textbf{Problem.} Outline a step-by-step build of a loan portfolio credit risk model.

\textbf{Solution.}

\begin{enumerate}
\item \textbf{Data.} Gather borrower financials, collateral information, vintage and performance data, and macroeconomic variables. Define default and cure consistently.

\item \textbf{PD model.} Build a logistic or survival (hazard) model with borrower and macro covariates; calibrate TTC versus PIT; validate with AUC, calibration, and stability tests.

\item \textbf{LGD model.} Use beta regression, GLMs, or quantile regression based on collateral, seniority, and LTV; constrain predictions to \((0,1)\); validate with pinball loss and calibration diagnostics.

\item \textbf{EAD / CCF.} For revolvers, model drawdown at default; for amortizing loans, project the exposure path.

\item \textbf{Portfolio aggregation.} Use factor or copula dependence structures to simulate joint defaults and compute expected loss, unexpected loss, full loss distributions, and concentration metrics.

\item \textbf{Stress testing.} Apply macro scenarios and report impacts on capital, provisions, and loss distributions.
\end{enumerate}

\begin{EditorialNote}


\textbf{Core skill tested.} This is a risk-and-finance framing problem. The best answer defines the object precisely, states the minimal formula or workflow, and then immediately mentions the practical assumption set, because interviewers care about whether you know where textbook simplifications break in a desk setting.

\textbf{How to present it in interview.} Answer these by sequencing definition -> construction -> caveat. For example, define the measure or sensitivity, say how you would compute it on real data or positions, and then mention one limitation such as liquidity, model error, or convention dependence.

\textbf{Common traps.}
\begin{itemize}
\item Quoting a definition without specifying horizon, confidence level, or market convention.
\item Ignoring model risk, liquidity, or scenario design when moving from textbook to practice.
\item Giving a sensitivity or valuation formula without saying what inputs must be bootstrapped or estimated.
\end{itemize}

\textbf{Related follow-ups.}
\begin{itemize}
\item What part of the workflow is convention-dependent, and what part is model-dependent?
\item How would you explain the same concept to a portfolio manager instead of a quant interviewer?
\end{itemize}


\end{EditorialNote}
\QuestionSpacer

\subsection{Fixed income, duration, and rates concepts}

\InterviewQuestionHeader{Question 22}{Global investment bank and financial services firm}{Risk / Quant}{Duration}{Easy}
\textbf{Problem.} Define duration variants and how to use them.

\textbf{Solution.}

Macaulay duration is the present-value-weighted average time to receive the bond's cash flows.

Modified duration is the first-order sensitivity of bond price to a parallel yield shift. Under \(m\)-times compounding,
\[
D_{\mathrm{mod}}=\frac{D_{\mathrm{Mac}}}{1+y/m}.
\]

DV01 is approximately
\[
\mathrm{DV01}\approx P \cdot D_{\mathrm{mod}} \cdot 0.0001.
\]

Key-rate durations isolate sensitivity to selected maturities on the yield curve and are useful for curve hedging.

In practice, duration is used for immunization, asset-liability matching, hedge design, and convexity-aware interest-rate risk management.

\begin{EditorialNote}


\textbf{Core skill tested.} This is a risk-and-finance framing problem. The best answer defines the object precisely, states the minimal formula or workflow, and then immediately mentions the practical assumption set, because interviewers care about whether you know where textbook simplifications break in a desk setting.

\textbf{How to present it in interview.} Answer these by sequencing definition -> construction -> caveat. For example, define the measure or sensitivity, say how you would compute it on real data or positions, and then mention one limitation such as liquidity, model error, or convention dependence.

\textbf{Common traps.}
\begin{itemize}
\item Quoting a definition without specifying horizon, confidence level, or market convention.
\item Ignoring model risk, liquidity, or scenario design when moving from textbook to practice.
\item Giving a sensitivity or valuation formula without saying what inputs must be bootstrapped or estimated.
\end{itemize}

\textbf{Related follow-ups.}
\begin{itemize}
\item What part of the workflow is convention-dependent, and what part is model-dependent?
\item How would you explain the same concept to a portfolio manager instead of a quant interviewer?
\end{itemize}


\end{EditorialNote}
\QuestionSpacer

\subsection{Portfolio construction, factor models, and attribution}

\InterviewQuestionHeader{Question 23}{Quantitative investment management firm}{Quantitative Researcher}{Portfolio Optimization}{Medium}
\textbf{Problem.} You have \(n\) assets with expected returns \(\mu \in \mathbb R^n\) and covariance \(\Sigma \succ 0\). Find the minimum-variance portfolio for a target expected return \(r\) (and full investment).

\textbf{Solution.}

Solve
\[
\min_w \; w^\top \Sigma w
\]
subject to
\[
\mu^\top w = r,
\qquad
\mathbf 1^\top w = 1.
\]

The Lagrangian is
\[
L(w,\lambda,\nu)
=
w^\top \Sigma w
-
\lambda(\mu^\top w-r)
-
\nu(\mathbf 1^\top w-1).
\]

The first-order condition is
\[
2\Sigma w - \lambda \mu - \nu \mathbf 1 = 0,
\]
so
\[
w = \frac12 \Sigma^{-1}(\lambda \mu + \nu \mathbf 1).
\]

Define
\[
A = \mathbf 1^\top \Sigma^{-1}\mathbf 1,
\qquad
B = \mathbf 1^\top \Sigma^{-1}\mu,
\qquad
C = \mu^\top \Sigma^{-1}\mu,
\qquad
\Delta = AC-B^2.
\]

Imposing the two constraints gives
\[
\lambda = \frac{2(Ar-B)}{\Delta},
\qquad
\nu = \frac{2(C-Br)}{\Delta}.
\]

Hence the optimal weights are
\[
w^\ast
=
\Sigma^{-1}
\left(
\frac{Ar-B}{\Delta}\mu
+
\frac{C-Br}{\Delta}\mathbf 1
\right).
\]

The minimum variance achieved at target return \(r\) is
\[
\sigma^2(r)
=
\frac{Ar^2-2Br+C}{\Delta}.
\]

\begin{EditorialNote}


\textbf{Core skill tested.} This is a risk-and-finance framing problem. The best answer defines the object precisely, states the minimal formula or workflow, and then immediately mentions the practical assumption set, because interviewers care about whether you know where textbook simplifications break in a desk setting.

\textbf{How to present it in interview.} Answer these by sequencing definition -> construction -> caveat. For example, define the measure or sensitivity, say how you would compute it on real data or positions, and then mention one limitation such as liquidity, model error, or convention dependence.

\textbf{Common traps.}
\begin{itemize}
\item Quoting a definition without specifying horizon, confidence level, or market convention.
\item Ignoring model risk, liquidity, or scenario design when moving from textbook to practice.
\item Giving a sensitivity or valuation formula without saying what inputs must be bootstrapped or estimated.
\end{itemize}

\textbf{Related follow-ups.}
\begin{itemize}
\item What part of the workflow is convention-dependent, and what part is model-dependent?
\item How would you explain the same concept to a portfolio manager instead of a quant interviewer?
\end{itemize}


\end{EditorialNote}
\QuestionSpacer

\section{Trading, Market Making, and Statistical Arbitrage}
\ChapterIntro{Use these questions to practice complete solutions, concise explanations, and likely follow-up questions.}

\subsection{Market making, hedging, and microstructure}

\InterviewQuestionHeader{Question 24}{Quantitative proprietary trading firm}{Quantitative Researcher}{Market Microstructure}{Medium}
\textbf{Problem.} Explain the concept of market microstructure and its impact on trading.

\textbf{Solution.}

\begin{itemize}
\item \textbf{Core elements:} limit-order book, tick size, queue priority (price-time), maker-taker fees, hidden or iceberg orders, and latency.

\item \textbf{Empirical effects:} bid-ask bounce (which induces negative first-order autocorrelation in transaction prices), price impact \(\Lambda\), the spread-volatility relation, and variance-ratio distortions at high frequency.

\item \textbf{Adverse selection and impact models:}
\begin{itemize}
\item Kyle model:
\[
\Delta P = \Lambda Q + \eta;
\]
\item Almgren--Chriss: temporary and permanent impact with optimal execution;
\item Hasbrouck-style trade-sign models for impact decay.
\end{itemize}

\item \textbf{Trading implications:} choose the execution style (limit order versus market order), manage participation rate and queue position, avoid toxic flow, and measure implementation shortfall relative to benchmarks such as open, close, or VWAP.

\item \textbf{Backtesting caution:} use midquotes or quote-consistent fills, and apply realistic queue, latency, and partial-fill rules. Microstructure determines costs, slippage, and execution risk, so modeling it accurately is essential for alpha preservation and capacity estimation.
\end{itemize}

\begin{EditorialNote}


\textbf{Core skill tested.} This prompt is less about memorizing a single formula and more about demonstrating trading judgment. You want to show that signal generation, execution costs, inventory control, and validation all live in the same answer rather than appearing as disconnected bullet points.

\textbf{How to present it in interview.} Treat this like a small investment memo: define the edge, state how you would implement it, then explain how you would keep yourself from being fooled by costs, stale signals, or biased backtests. The interviewer is usually looking for disciplined judgment rather than a single clever sentence.

\textbf{Common traps.}
\begin{itemize}
\item Talking about signal quality without discussing transaction costs, latency, or inventory risk.
\item Describing a backtest that implicitly leaks future information or overfits thresholds.
\item Presenting a strategy idea without a falsification plan or risk control.
\end{itemize}

\textbf{Related follow-ups.}
\begin{itemize}
\item What is the falsification test that would make you stop trusting the signal or quoting rule?
\item How would you adapt the idea if spreads widen, liquidity disappears, or the regime changes?
\end{itemize}


\end{EditorialNote}
\QuestionSpacer

\InterviewQuestionHeader{Question 25}{Quantitative proprietary trading and market-making firm}{Quantitative Trader}{Market Making / Bayesian Pricing}{Medium}
\textbf{Problem.} Market making scenario: A stock's "true value" \(V\) is modeled as \(V \sim N(\mu,\sigma^2)\) with \(\mu=100\), \(\sigma=2\). You must quote bid/ask prices. What would you quote?

\textbf{Solution.}


The prior \(V\sim N(100,2^2)\) determines a natural \emph{center} for the quotes, but it does not determine a unique bid--ask spread.

If the market maker is risk-neutral, has no inventory, faces no fees or latency, and incoming orders are pure noise independent of \(V\), competition drives the zero-profit quotes to
\[
\text{bid}=\text{ask}=\mathbb E[V]=100.
\]
A positive spread requires an additional economic friction or model: adverse selection, inventory risk, order-processing cost, tick size, latency risk, or a fill-intensity trade-off.

For example, in an adverse-selection model the competitive zero-profit quotes take the form
\[
\text{ask}=\mathbb E[V\mid\text{incoming buy}],
\qquad
\text{bid}=\mathbb E[V\mid\text{incoming sell}],
\]
and these conditional expectations cannot be computed until \(P(\text{order side}\mid V)\) is specified.

Therefore the interview-quality answer is:
\begin{itemize}
\item anchor the mid at \(100\) under the symmetric prior;
\item state that the spread is \emph{not identified} from \(\mu\) and \(\sigma\) alone;
\item ask for, or explicitly assume, a fill/adverse-selection/cost model before quoting a numerical bid and ask.
\end{itemize}
A confidence interval such as \(100\pm z\sigma\) may be useful as a risk diagnostic, but it is not by itself a zero-profit market-making quote.

\begin{EditorialNote}


\textbf{Core skill tested.} This prompt is less about memorizing a single formula and more about demonstrating trading judgment. You want to show that signal generation, execution costs, inventory control, and validation all live in the same answer rather than appearing as disconnected bullet points.

\textbf{How to present it in interview.} Treat this like a small investment memo: define the edge, state how you would implement it, then explain how you would keep yourself from being fooled by costs, stale signals, or biased backtests. The interviewer is usually looking for disciplined judgment rather than a single clever sentence.

\textbf{Common traps.}
\begin{itemize}
\item Talking about signal quality without discussing transaction costs, latency, or inventory risk.
\item Describing a backtest that implicitly leaks future information or overfits thresholds.
\item Presenting a strategy idea without a falsification plan or risk control.
\end{itemize}

\textbf{Related follow-ups.}
\begin{itemize}
\item What is the falsification test that would make you stop trusting the signal or quoting rule?
\item How would you adapt the idea if spreads widen, liquidity disappears, or the regime changes?
\end{itemize}


\end{EditorialNote}
\QuestionSpacer

\subsection{Statistical arbitrage, signals, and backtesting}

\InterviewQuestionHeader{Question 26}{Quantitative proprietary trading firm}{Quantitative Researcher}{Backtesting / Experimental Design}{Medium}
\textbf{Problem.} How do you properly backtest a trading strategy?

\textbf{Solution.}

\begin{enumerate}
\item \textbf{Data hygiene.} Use point-in-time data; handle survivorship bias, delistings, stale quotes, corporate actions, and dividends. Eliminate look-ahead bias and leakage.

\item \textbf{Costs and execution.} Include realistic transaction costs, commissions, fees, slippage, borrow and financing costs, shorting constraints, and an execution model with market impact.

\item \textbf{Sampling and latency.} Respect exchange timestamps, signal-construction delays, order-placement latency, and executable prices rather than idealized prints.

\item \textbf{Evaluation design.} Use train/validation/test splits or walk-forward evaluation. Freeze hyperparameters before the test period and report out-of-sample performance only.

\item \textbf{Robustness.} Perform parameter sweeps, perturb costs and latency assumptions, break results down by subperiod and regime, and use bootstrap confidence intervals or reality checks.

\item \textbf{Multiple testing control.} Control data snooping using White's Reality Check, SPA tests, FDR control, or clean holdout procedures.

\item \textbf{Risk and capacity.} Monitor drawdowns, turnover, exposure limits, liquidity filters such as ADV or participation caps, concentration, and borrow availability.

\item \textbf{Reporting.} Report metrics such as IC, IR, Sharpe, Sortino, maximum drawdown, hit rate, skew, kurtosis, and \(t\)-statistics with Newey--West corrections for autocorrelation and heteroskedasticity.
\end{enumerate}

\begin{EditorialNote}


\textbf{Core skill tested.} This prompt is less about memorizing a single formula and more about demonstrating trading judgment. You want to show that signal generation, execution costs, inventory control, and validation all live in the same answer rather than appearing as disconnected bullet points.

\textbf{How to present it in interview.} Treat this like a small investment memo: define the edge, state how you would implement it, then explain how you would keep yourself from being fooled by costs, stale signals, or biased backtests. The interviewer is usually looking for disciplined judgment rather than a single clever sentence.

\textbf{Common traps.}
\begin{itemize}
\item Talking about signal quality without discussing transaction costs, latency, or inventory risk.
\item Describing a backtest that implicitly leaks future information or overfits thresholds.
\item Presenting a strategy idea without a falsification plan or risk control.
\end{itemize}

\textbf{Related follow-ups.}
\begin{itemize}
\item What is the falsification test that would make you stop trusting the signal or quoting rule?
\item How would you adapt the idea if spreads widen, liquidity disappears, or the regime changes?
\end{itemize}


\end{EditorialNote}
\QuestionSpacer

\InterviewQuestionHeader{Question 27}{Quantitative investment management firm}{Quantitative Researcher}{Backtesting / Evaluation Metrics}{Medium}
\textbf{Problem.} Describe how you would backtest a trading strategy and which metrics you would use to evaluate performance.

\textbf{Solution.}

\textbf{Protocol.}

\begin{enumerate}
\item \textbf{Data hygiene.} Use survivorship-bias-free universes, point-in-time fundamentals, correct timestamp alignment, corporate-action adjustments, and realistic borrow or financing availability.

\item \textbf{Split design.} Use expanding or rolling walk-forward splits:
\[
\text{train} \rightarrow \text{validate} \rightarrow \text{test}.
\]
Fit all feature engineering and preprocessing only on in-sample data to avoid leakage.

\item \textbf{Execution model.} Include realistic commissions, fees, spreads, slippage, latency, market impact, borrow fees, funding costs, and shorting constraints.

\item \textbf{Portfolio rules.} Define sizing, leverage, turnover constraints, concentration limits, exposure caps, and rebalancing schedules.

\item \textbf{Risk control.} Add pre-trade limits, post-trade kill switches, scenario tests, and regime segmentation.
\end{enumerate}

\textbf{Key metrics.}

\begin{itemize}
\item \textbf{Performance:} CAGR, Sharpe, Sortino, information ratio, CAPM alpha and beta, hit rate, average win/loss, and profit factor.
\item \textbf{Risk:} volatility, maximum drawdown, Calmar ratio, VaR, Expected Shortfall, and exposures by sector, country, or factor.
\item \textbf{Stability and robustness:} rolling Sharpe, turnover, capacity versus impact, parameter sensitivity, and P\&L by time, volatility, or liquidity bucket.
\item \textbf{Multiple testing control:} FDR or FWER control across many ideas, together with live-versus-backtest decay tracking.
\end{itemize}

\begin{EditorialNote}


\textbf{Core skill tested.} This prompt is less about memorizing a single formula and more about demonstrating trading judgment. You want to show that signal generation, execution costs, inventory control, and validation all live in the same answer rather than appearing as disconnected bullet points.

\textbf{How to present it in interview.} Treat this like a small investment memo: define the edge, state how you would implement it, then explain how you would keep yourself from being fooled by costs, stale signals, or biased backtests. The interviewer is usually looking for disciplined judgment rather than a single clever sentence.

\textbf{Common traps.}
\begin{itemize}
\item Talking about signal quality without discussing transaction costs, latency, or inventory risk.
\item Describing a backtest that implicitly leaks future information or overfits thresholds.
\item Presenting a strategy idea without a falsification plan or risk control.
\end{itemize}

\textbf{Related follow-ups.}
\begin{itemize}
\item What is the falsification test that would make you stop trusting the signal or quoting rule?
\item How would you adapt the idea if spreads widen, liquidity disappears, or the regime changes?
\end{itemize}


\end{EditorialNote}
\QuestionSpacer

\section{Algorithms, Data Structures, and Interview Coding}
\ChapterIntro{Use these questions to practice complete solutions, concise explanations, and likely follow-up questions.}

\subsection{Algorithms, data structures, and coding patterns}

\InterviewQuestionHeader{Question 28}{Options market-making and proprietary trading firm}{Quantitative Trader}{Data Structures / Streaming Median}{Medium}
\textbf{Problem.} Given a stream of numbers, design a data structure that can return the median after each insertion, efficiently.

\textbf{Solution.}

Maintain two heaps:
\begin{itemize}
\item a max-heap \(L\) for the lower half of the data,
\item a min-heap \(R\) for the upper half.
\end{itemize}

Keep the invariants
\[
|\mathrm{len}(L)-\mathrm{len}(R)| \le 1,
\qquad
\max(L) \le \min(R).
\]

On insertion:
\begin{enumerate}
\item insert the new value into the appropriate heap;
\item rebalance if one heap exceeds the other by \(2\) elements.
\end{enumerate}

The median is then:
\[
\max(L)
\quad \text{if } L \text{ has one extra element,}
\]
\[
\min(R)
\quad \text{if } R \text{ has one extra element,}
\]
and
\[
\frac{\max(L)+\min(R)}{2}
\quad \text{if the heaps have equal size.}
\]

Complexities:
\[
O(\log n)
\quad \text{per insertion,}
\qquad
O(1)
\quad \text{per median query,}
\qquad
O(n)
\quad \text{space.}
\]

\begin{EditorialNote}


\textbf{Core skill tested.} This is really an invariants-and-complexity question disguised as implementation. The interviewer wants to hear what state you maintain, why that state is sufficient, and how each operation preserves the invariant while meeting the target time complexity.

\textbf{How to present it in interview.} Present these problems in three moves: name the invariant, describe insert/update/query, and give time and space complexity. If there is a balancing rule or a lazy-deletion rule, say it before discussing edge cases so the interviewer can see the design is principled rather than improvised.

\textbf{Common traps.}
\begin{itemize}
\item Explaining the container choice but not the invariant that makes the query correct.
\item Forgetting duplicates, empty-state behaviour, or rebalance conditions.
\item Omitting time and space complexity or misstating amortized versus worst-case cost.
\end{itemize}

\textbf{Related follow-ups.}
\begin{itemize}
\item Can you support deletion, duplicates, or sliding-window updates without breaking the complexity target?
\item How would the design change if the stream were distributed across shards or needed persistence?
\end{itemize}


\end{EditorialNote}
\QuestionSpacer

\InterviewQuestionHeader{Question 29}{Quantitative proprietary trading firm}{Quantitative Researcher}{Algorithms / DP}{Medium}
\textbf{Problem.} Given an integer array, find the maximum sum over all contiguous subarrays (Kadane's algorithm). Now modify it for the case where you may remove at most one element from the chosen subarray. Return the maximum attainable sum.

\textbf{Solution.}

Maintain two DP states at position \(i\), both for subarrays ending at \(i\):
\[
\mathrm{noDel}[i]
:= \text{maximum subarray sum ending at } i \text{ with no deletion},
\]
\[
\mathrm{oneDel}[i]
:= \text{maximum subarray sum ending at } i \text{ with at most one deletion}.
\]

For \(i \ge 1\), letting \(x=a_i\), the transitions are
\[
\mathrm{noDel}[i] = \max\{x,\ \mathrm{noDel}[i-1]+x\},
\]
\[
\mathrm{oneDel}[i] = \max\{\mathrm{noDel}[i-1],\ \mathrm{oneDel}[i-1]+x\}.
\]

Interpretation:
\begin{itemize}
\item For \(\mathrm{noDel}\), either start a new subarray at \(x\) or extend the previous one.
\item For \(\mathrm{oneDel}\), either delete the current element \(x\) and keep \(\mathrm{noDel}[i-1]\), or keep \(x\) and extend a subarray that has already used one deletion.
\end{itemize}

Initialization:
\[
\mathrm{noDel}[0] = a_0,
\qquad
\mathrm{oneDel}[0] = -\infty.
\]

The answer is
\[
\max_i \max\{\mathrm{noDel}[i],\ \mathrm{oneDel}[i]\}.
\]

This runs in
\[
O(n)
\]
time and
\[
O(1)
\]
space if only rolling values are stored.

Reference implementation (Python, \(O(n)\) time, \(O(1)\) space):
\begin{codeblock}
def max_subarray_with_one_deletion(a):
    no_del = a[0]
    one_del = float('-inf')
    ans = a[0]

    for x in a[1:]:
        one_del = max(no_del, one_del + x)   # delete x OR keep x after prior deletion
        no_del = max(x, no_del + x)          # classic Kadane step
        ans = max(ans, no_del, one_del)

    return ans
\end{codeblock}

\begin{EditorialNote}


\textbf{Core skill tested.} This is really an invariants-and-complexity question disguised as implementation. The interviewer wants to hear what state you maintain, why that state is sufficient, and how each operation preserves the invariant while meeting the target time complexity.

\textbf{How to present it in interview.} Present these problems in three moves: name the invariant, describe insert/update/query, and give time and space complexity. If there is a balancing rule or a lazy-deletion rule, say it before discussing edge cases so the interviewer can see the design is principled rather than improvised.

\textbf{Common traps.}
\begin{itemize}
\item Explaining the container choice but not the invariant that makes the query correct.
\item Forgetting duplicates, empty-state behaviour, or rebalance conditions.
\item Omitting time and space complexity or misstating amortized versus worst-case cost.
\end{itemize}

\textbf{Related follow-ups.}
\begin{itemize}
\item Can you support deletion, duplicates, or sliding-window updates without breaking the complexity target?
\item How would the design change if the stream were distributed across shards or needed persistence?
\end{itemize}


\end{EditorialNote}
\QuestionSpacer

\InterviewQuestionHeader{Question 30}{Cross-industry quantitative-finance employers}{Quant Research}{Algorithms / Hashing}{Medium}
\textbf{Problem.} Given a list of words, group anagrams together (one line per group). Design and justify an \(O(nL)\) method where \(L\) is average word length.

\textbf{Solution.}

Two strings are anagrams if and only if they share the same canonical signature.

There are two common choices:

\begin{itemize}
\item \textbf{Sorted signature:} sort the characters of each word. This costs
\[
O(L\log L)
\]
per word.

\item \textbf{Count signature (preferred):} count characters into a fixed-length vector (for example, \(26\) entries for lowercase \(a\)-\(z\)) and use the resulting count tuple as the hash key. This costs
\[
O(L)
\]
per word.
\end{itemize}

Using the count signature, build a hash map
\[
\text{signature} \;\mapsto\; \text{list of words},
\]
which gives total complexity
\[
O(nL).
\]

Reference implementation (Python, count-key, case-insensitive):
\begin{codeblock}
from collections import defaultdict

def group_anagrams(words):
    buckets = defaultdict(list)
    for w in words:
        cnt = [0]*26
        for ch in w.lower():
            if 'a' <= ch <= 'z':
                cnt[ord(ch)-97] += 1
        buckets[tuple(cnt)].append(w)
    return list(buckets.values())

# Example:
# groups = group_anagrams(["earth","heart","ab","cat","listen","silent"])
# for g in groups: print(" ".join(g))
\end{codeblock}

\begin{EditorialNote}


\textbf{Core skill tested.} This is really an invariants-and-complexity question disguised as implementation. The interviewer wants to hear what state you maintain, why that state is sufficient, and how each operation preserves the invariant while meeting the target time complexity.

\textbf{How to present it in interview.} Present these problems in three moves: name the invariant, describe insert/update/query, and give time and space complexity. If there is a balancing rule or a lazy-deletion rule, say it before discussing edge cases so the interviewer can see the design is principled rather than improvised.

\textbf{Common traps.}
\begin{itemize}
\item Explaining the container choice but not the invariant that makes the query correct.
\item Forgetting duplicates, empty-state behaviour, or rebalance conditions.
\item Omitting time and space complexity or misstating amortized versus worst-case cost.
\end{itemize}

\textbf{Related follow-ups.}
\begin{itemize}
\item Can you support deletion, duplicates, or sliding-window updates without breaking the complexity target?
\item How would the design change if the stream were distributed across shards or needed persistence?
\end{itemize}


\end{EditorialNote}
\QuestionSpacer

\InterviewQuestionHeader{Question 31}{Quantitative investment management firm}{Quantitative Researcher}{Algorithms / Sliding Window}{Easy}
\textbf{Problem.} Implement a function to compute the moving average of a stream of integers given a window size \(k\).

\textbf{Solution.}

Maintain:
\begin{itemize}
\item a queue of at most \(k\) items,
\item a running sum of the values currently in the queue.
\end{itemize}

When a new value arrives:
\begin{enumerate}
\item push the new value into the queue;
\item add it to the running sum;
\item if the queue length exceeds \(k\), pop the oldest value and subtract it from the sum.
\end{enumerate}

The current moving average is then
\[
\frac{\text{sum}}{\text{len(queue)}}.
\]

Reference implementation:
\begin{codeblock}
from collections import deque

class MovingAverage:
    def __init__(self, k: int):
        self.k = k
        self.q = deque()
        self.s = 0.0

    def add(self, x: float) -> float:
        self.q.append(x)
        self.s += x
        if len(self.q) > self.k:
            self.s -= self.q.popleft()
        return self.s / len(self.q)
\end{codeblock}

Time per update:
\[
O(1),
\]
space:
\[
O(k).
\]

\begin{EditorialNote}


\textbf{Core skill tested.} This is really an invariants-and-complexity question disguised as implementation. The interviewer wants to hear what state you maintain, why that state is sufficient, and how each operation preserves the invariant while meeting the target time complexity.

\textbf{How to present it in interview.} Present these problems in three moves: name the invariant, describe insert/update/query, and give time and space complexity. If there is a balancing rule or a lazy-deletion rule, say it before discussing edge cases so the interviewer can see the design is principled rather than improvised.

\textbf{Common traps.}
\begin{itemize}
\item Explaining the container choice but not the invariant that makes the query correct.
\item Forgetting duplicates, empty-state behaviour, or rebalance conditions.
\item Omitting time and space complexity or misstating amortized versus worst-case cost.
\end{itemize}

\textbf{Related follow-ups.}
\begin{itemize}
\item Can you support deletion, duplicates, or sliding-window updates without breaking the complexity target?
\item How would the design change if the stream were distributed across shards or needed persistence?
\end{itemize}


\end{EditorialNote}
\QuestionSpacer

\section{Systems, Order Books, and Production Prototypes}
\ChapterIntro{Use these questions to practice complete solutions, concise explanations, and likely follow-up questions.}

\subsection{Real-time risk, OMS, and system architecture}

\InterviewQuestionHeader{Question 32}{Electronic market-making firm}{Quantitative Analyst}{Systems Design / Real-Time Risk}{Hard}
\textbf{Problem.} Design a real-time risk monitoring system that tracks positions across 200+ instruments and triggers alerts when limits are breached. The system must handle \textasciitilde{} 106 position updates per second.

\textbf{Solution.}


\textbf{Ingestion.}
\begin{itemize}
\item Multiple feed handlers (per venue/desk) parse updates
\[
\langle ts,\ inst\_id,\ acct,\ \Delta q,\ p,\ \gamma\text{-vector}\rangle.
\]
\item Each handler writes to an SPSC ring buffer (lock-free) consumed by the Risk Engine.
\item Sequence numbers provide idempotency, and duplicates are dropped.
\end{itemize}

\textbf{State (per instrument, optionally per account).}
\[
q_i \quad \text{current position (shares)},
\]
\[
n_i = p_i q_i \quad \text{notional},
\]
\[
\mathrm{Greeks}_i \approx J_i q_i \quad \text{linearized risk (e.g. DV01, Vega)}.
\]

All arrays are preallocated and keyed by dense
\[
inst\_id \in \{0,\dots,M-1\}
\]
to allow \(O(1)\) updates.

\textbf{Limits store.} Thresholds per instrument/account:
\[
|q_i| \le Q_i^{\max},
\qquad
|n_i| \le N_i^{\max},
\qquad
|DV01_i| \le D_i^{\max},
\qquad
|Vega_i| \le V_i^{\max}.
\]

Use hysteresis \((\pm \varepsilon)\) and debounce windows to avoid alert chatter.

\textbf{Risk Engine.}
\begin{itemize}
\item Reads batches from the rings, applies \(\Delta q\) to arrays, recomputes cheap linear metrics, and compares them to limits.
\item Breaches are pushed to an Alert Queue (SPSC), consumed by an out-of-band Notifier (email/chat/monitoring protocol).
\item Snapshots (every \(100\) to \(250\) ms) are published via shared memory/IPC for dashboards.
\item A WAL (write-ahead log) provides recovery.
\end{itemize}

\textbf{Throughput and latency targets.}
\begin{itemize}
\item Per update: \(O(1)\) arithmetic on a few cache-resident arrays, with no dynamic allocation.
\item Pin threads, isolate cores, use huge pages, and apply NUMA-aware placement.
\item Batch roughly
\[
N \approx 256
\]
records per read to amortize overhead.
\item For backpressure, rings expose occupancy; when near full, throttle feeders or drop to a degraded aggregate-only mode.
\end{itemize}

\textbf{Reference prototype (Python, hot path).} Array-based state; easily JIT-compiled with numba or ported to C++ for production.
\begin{codeblock}
import numpy as np
from dataclasses import dataclass

@dataclass
class Limits:
    q_max: np.ndarray
    notl_max: np.ndarray
    dv01_max: np.ndarray
    vega_max: np.ndarray
    # shape (M,)

class RealTimeRisk:
    def __init__(self, M: int, dv01_per_share, vega_per_share):
        self.M = M

        # state arrays
        self.q = np.zeros(M, dtype=np.float64)       # position
        self.price = np.zeros(M, dtype=np.float64)   # last price
        self.notl = np.zeros(M, dtype=np.float64)    # notional

        # linear Greek per share
        self.dv01_sh = np.asarray(dv01_per_share, dtype=np.float64)
        self.vega_sh = np.asarray(vega_per_share, dtype=np.float64)
        self.dv01 = np.zeros(M, dtype=np.float64)
        self.vega = np.zeros(M, dtype=np.float64)

        # limits (set later)
        self.lim = None

    def set_limits(self, lim: Limits):
        self.lim = lim

    # hot-path update: O(1)
    def on_update(self, inst_id: int, dq: float, price: float) -> int:
        i = inst_id
        self.q[i] += dq
        self.price[i] = price
        self.notl[i] = self.q[i] * price
        self.dv01[i] = self.q[i] * self.dv01_sh[i]
        self.vega[i] = self.q[i] * self.vega_sh[i]
        return self._check(i)

    # returns bitmask of breaches for instrument i
    def _check(self, i: int) -> int:
        if self.lim is None:
            return 0
        mask = 0
        eps = 1e-9
        if abs(self.q[i]) > self.lim.q_max[i] + eps:
            mask |= 1 << 0
        if abs(self.notl[i]) > self.lim.notl_max[i] + eps:
            mask |= 1 << 1
        if abs(self.dv01[i]) > self.lim.dv01_max[i] + eps:
            mask |= 1 << 2
        if abs(self.vega[i]) > self.lim.vega_max[i] + eps:
            mask |= 1 << 3
        return mask
\end{codeblock}

\textbf{Alerting \& reliability.}
\begin{itemize}
\item Each breach produces
\[
\langle ts,\ inst\_id,\ metric,\ val,\ limit\rangle
\]
on an Alert Queue; the notifier applies hysteresis/debounce and routes by desk.
\item Recovery: append every applied update to a WAL; on restart, load the last snapshot and replay the WAL to reconstruct state (typically \(\le 1\) s).
\item Monitoring: export per-stage latency, ring occupancy, and breach rates; auto-escalate when the risk engine lags.
\end{itemize}

\textbf{Why it scales to \(10^6\)/s.}
\begin{itemize}
\item Updates touch only a constant number of cache-resident doubles (positions, two Greeks, and notional).
\item Batching and SPSC rings remove locks.
\item Preallocated arrays avoid GC/allocation overhead.
\item Horizontal scaling is achieved by sharding instruments across cores; global aggregates can be computed by a second-stage reducer at lower frequency.
\end{itemize}

\begin{EditorialNote}


\textbf{Core skill tested.} This question tests whether you can turn a production-style requirement into a state model with explicit invariants. A strong answer names the entities you store, the updates you apply on each event, and the correctness conditions you refuse to violate even under latency pressure.

\textbf{How to present it in interview.} In a live interview, open with the data model, then walk through the event path in order: ingest, normalize, update, publish, and recover from gaps or failures. Say the words correctness, sequencing, and recovery explicitly; that is usually what differentiates a trading-systems answer from a toy coding answer.

\textbf{Common traps.}
\begin{itemize}
\item Giving a fast design without specifying what state is authoritative under reorders, cancels, or sequence gaps.
\item Talking only about happy-path latency and ignoring recovery, backpressure, and stale-state publication.
\item Assuming the sketch is production-safe without naming invariants, clocking, or correctness checks.
\end{itemize}

\textbf{Related follow-ups.}
\begin{itemize}
\item What changes if messages arrive out of order or you need to recover from a dropped sequence range?
\item How would you test correctness and latency without relying on a single happy-path benchmark?
\end{itemize}


\end{EditorialNote}
\QuestionSpacer

\subsection{Order books, VWAP, and production market-data infrastructure}

\InterviewQuestionHeader{Question 33}{Quantitative investment management firm}{Quantitative Researcher}{Systems / Order Book}{Medium}
\textbf{Problem.} Implement a limit order book class supporting add, cancel, and match with price-time priority.

\textbf{Solution.}


Use FIFO queues at each price level, heaps for discovering the best bid/ask, and an order-id map that points directly to mutable order objects. Cancellation can then mark an order dead in expected \(O(1)\), while matching purges canceled orders from the front before executing.

\begin{codeblock}
from collections import defaultdict, deque
import heapq

class LOB:
    def __init__(self):
        self.bids = defaultdict(deque)
        self.asks = defaultdict(deque)
        self.bid_heap = []          # negative prices
        self.ask_heap = []
        self.orders = {}            # id -> [id, qty, timestamp, side, price]
        self.clock = 0

    def add(self, oid, side, price, qty):
        if oid in self.orders or qty <= 0:
            raise ValueError("invalid order")
        order = [oid, qty, self.clock, side, price]
        self.clock += 1
        self.orders[oid] = order
        book = self.bids if side == 'B' else self.asks
        book[price].append(order)
        heapq.heappush(self.bid_heap if side == 'B' else self.ask_heap,
                       -price if side == 'B' else price)
        return self.match()

    def cancel(self, oid):
        order = self.orders.pop(oid, None)
        if order is not None:
            order[1] = 0            # lazy deletion; O(1)

    def _best(self, side):
        book = self.bids if side == 'B' else self.asks
        heap = self.bid_heap if side == 'B' else self.ask_heap
        while heap:
            price = -heap[0] if side == 'B' else heap[0]
            q = book.get(price)
            while q and q[0][1] == 0:
                q.popleft()
            if q:
                return price
            heapq.heappop(heap)
            book.pop(price, None)
        return None

    def match(self):
        fills = []
        while True:
            bid = self._best('B')
            ask = self._best('S')
            if bid is None or ask is None or bid < ask:
                break
            b = self.bids[bid][0]
            a = self.asks[ask][0]
            qty = min(b[1], a[1])
            # The older order is resting; trade at its limit price.
            trade_price = bid if b[2] < a[2] else ask
            fills.append((b[0], a[0], trade_price, qty))
            b[1] -= qty
            a[1] -= qty
            if b[1] == 0:
                self.orders.pop(b[0], None)
                self.bids[bid].popleft()
            if a[1] == 0:
                self.orders.pop(a[0], None)
                self.asks[ask].popleft()
        return fills
\end{codeblock}

Price-time priority is enforced by selecting the best price first and then the FIFO order at that price. Heap entries are lazily invalidated. In a production exchange simulator, also define modify semantics, self-trade prevention, tick-size validation, partial cancels, and deterministic sequencing.

\begin{EditorialNote}


\textbf{Core skill tested.} This question tests whether you can turn a production-style requirement into a state model with explicit invariants. A strong answer names the entities you store, the updates you apply on each event, and the correctness conditions you refuse to violate even under latency pressure.

\textbf{How to present it in interview.} In a live interview, open with the data model, then walk through the event path in order: ingest, normalize, update, publish, and recover from gaps or failures. Say the words correctness, sequencing, and recovery explicitly; that is usually what differentiates a trading-systems answer from a toy coding answer.

\textbf{Common traps.}
\begin{itemize}
\item Giving a fast design without specifying what state is authoritative under reorders, cancels, or sequence gaps.
\item Talking only about happy-path latency and ignoring recovery, backpressure, and stale-state publication.
\item Assuming the sketch is production-safe without naming invariants, clocking, or correctness checks.
\end{itemize}

\textbf{Related follow-ups.}
\begin{itemize}
\item What changes if messages arrive out of order or you need to recover from a dropped sequence range?
\item How would you test correctness and latency without relying on a single happy-path benchmark?
\end{itemize}


\end{EditorialNote}
\QuestionSpacer

\section{Professional Skills, Governance, and Communication}
\ChapterIntro{Use these questions to practice complete solutions, concise explanations, and likely follow-up questions.}

\subsection{Behavioral, communication, and research process}

\InterviewQuestionHeader{Question 34}{Global investment bank and financial services firm}{Risk / Quant}{Behavioural}{Easy}
\textbf{Problem.} Describe a time you identified a significant risk others missed (STAR format).

\textbf{Solution.}

\begin{enumerate}
\item \textbf{Situation:} context and why it mattered.
\item \textbf{Task:} your responsibility and target outcome.
\item \textbf{Action:} data you pulled, checks you ran, stakeholders engaged, and what you built (e.g. stress/monitor).
\item \textbf{Result:} quantified impact avoided, changes adopted, and what you learned.
\end{enumerate}

\begin{EditorialNote}


\textbf{Core skill tested.} This is a risk-and-finance framing problem. The best answer defines the object precisely, states the minimal formula or workflow, and then immediately mentions the practical assumption set, because interviewers care about whether you know where textbook simplifications break in a desk setting.

\textbf{How to present it in interview.} Answer these by sequencing definition -> construction -> caveat. For example, define the measure or sensitivity, say how you would compute it on real data or positions, and then mention one limitation such as liquidity, model error, or convention dependence.

\textbf{Common traps.}
\begin{itemize}
\item Quoting a definition without specifying horizon, confidence level, or market convention.
\item Ignoring model risk, liquidity, or scenario design when moving from textbook to practice.
\item Giving a sensitivity or valuation formula without saying what inputs must be bootstrapped or estimated.
\end{itemize}

\textbf{Related follow-ups.}
\begin{itemize}
\item What part of the workflow is convention-dependent, and what part is model-dependent?
\item How would you explain the same concept to a portfolio manager instead of a quant interviewer?
\end{itemize}


\end{EditorialNote}

\QuestionSpacer

\subsection{Model validation, regulation, and governance}

\InterviewQuestionHeader{Question 35}{Global investment bank and financial services firm}{Risk / Quant}{Model Validation \& Governance}{Medium}
\textbf{Problem.} Describe your approach to model validation and governance.

\textbf{Solution.}

\begin{enumerate}
\item \textbf{Conceptual soundness:} theory, assumptions, intended use, and model boundaries.
\item \textbf{Data quality:} lineage, representativeness, missing values and outliers, point-in-time handling, and survivorship-bias controls.
\item \textbf{Implementation verification:} code review, benchmark tests, and numerical-stability checks.
\item \textbf{Outcomes analysis:} backtests, performance metrics, stability or PSI, and stress behavior.
\item \textbf{Ongoing monitoring:} KPIs, KRIs, drift detection, and periodic recalibration triggers.
\item \textbf{Governance:} inventories, documentation, change control, independent review, and periodic re-approval.
\end{enumerate}

\begin{EditorialNote}


\textbf{Core skill tested.} This item tests whether you can move cleanly from definition to estimation to validation. A polished response states the estimator or framework, clarifies its assumptions, and then points out the most likely failure mode such as leakage, non-stationarity, or unstable inference.

\textbf{How to present it in interview.} A concise interview answer should always separate modelling from evaluation. State the model or diagnostic, then immediately describe the validation design and the reason you trust or distrust the result under realistic data conditions.

\textbf{Common traps.}
\begin{itemize}
\item Stating an estimator but not the assumptions needed for interpretation.
\item Forgetting to separate train/validation/test logic or time ordering.
\item Choosing a metric that does not match the business or modelling objective.
\end{itemize}

\textbf{Related follow-ups.}
\begin{itemize}
\item What breaks first if the data are non-stationary, heavy-tailed, or missing not at random?
\item How would you redesign the validation scheme if the labels overlap or the data are strictly time ordered?
\end{itemize}


\end{EditorialNote}
\QuestionSpacer
\section*{Cross-Industry Foundation Questions}
\addcontentsline{toc}{section}{Cross-Industry Foundation Questions}
The final 15 questions reinforce mathematical and financial ideas that recur across institution types. They are deliberately chosen for breadth rather than to reproduce every variation of a familiar puzzle or derivation.
\section{Brain Teasers}

\ChapterIntro{Use these questions to practice complete solutions, concise explanations, and likely follow-up questions.}

\subsection{Problem Simplification}

\FoundationQuestionHeader{Question 36}{Screwy Pirates}{Problem Simplification}

\textbf{Problem.} Five perfectly rational pirates must divide 100 gold coins. The most senior pirate proposes a split and all pirates vote, including the proposer. A proposal passes with at least 50\% of the votes; otherwise the proposer is killed and the process repeats with the next pirate. Preferences are lexicographic: survive first, maximize gold second, and, if otherwise indifferent, prefer fewer surviving pirates. How are the 100 coins divided among five pirates?


\textbf{Solution.}

\textbf{Answer.} The most senior pirate keeps 98 coins and gives one coin each to the two pirates who would receive nothing if he were killed; the remaining two pirates get zero.

\textbf{Reasoning.} Work backward. With two pirates, the senior pirate can approve his own proposal. With three pirates, the new senior pirate buys the vote of the pirate who would get zero in the two-pirate outcome for one coin. The same logic alternates as pirates are added. In the five-pirate case, the senior pirate needs three votes including his own, so he buys the two cheapest votes for one coin each and keeps the remaining 98.


\begin{EditorialNote}


\textbf{Core skill tested.} This problem tests structural reasoning: identify the invariant, symmetry, backward-induction rule, or information constraint before doing arithmetic.


\textbf{How to present it in interview.} State the key structural observation first, reduce to a smaller or symmetric case if possible, then finish with a short proof or construction.


\textbf{Common traps.}
\begin{itemize}
\item Brute-forcing cases before identifying the invariant.
\item Assuming symmetry that the rules do not actually provide.
\item Giving a construction without proving that it is minimal or guaranteed.
\end{itemize}


\textbf{Related follow-ups.}
\begin{itemize}
\item Can you generalize the argument to n objects or a different threshold?
\item Is there a second solution using an invariant, parity, or information-theoretic bound?
\end{itemize}


\end{EditorialNote}
\QuestionSpacer

\subsection{Logic Reasoning}

\FoundationQuestionHeader{Question 37}{Defective Ball}{Logic Reasoning}

\textbf{Problem.} Among 12 visually identical balls, exactly one has a different weight and may be either heavier or lighter. Using only a balance scale, identify the odd ball and whether it is heavy or light in at most three weighings.


\textbf{Solution.}

Label the balls 1--12. First weigh \(1,2,3,4\) against \(5,6,7,8\).

If they balance, the odd ball is among 9--12. Weigh \(9,10,11\) against three known-normal balls. If that balances, 12 is odd and a final comparison with a normal ball identifies whether it is heavy or light. If it does not balance, the direction tells whether 9--11 are heavy or light; compare two of them to identify the ball.

If the first weighing is left-heavy, the candidates are \(1,2,3,4\) heavy or \(5,6,7,8\) light. Weigh \(1,2,5\) against \(3,6,9\), where 9 is known normal. A left-heavy result leaves \(1_H,2_H,6_L\); a right-heavy result leaves \(3_H,5_L\); a balance leaves \(4_H,7_L,8_L\). One final comparison resolves each case. The right-heavy first-weighing case is symmetric. Thus three weighings always suffice.


\begin{EditorialNote}


\textbf{Core skill tested.} This problem tests structural reasoning: identify the invariant, symmetry, backward-induction rule, or information constraint before doing arithmetic.


\textbf{How to present it in interview.} State the key structural observation first, reduce to a smaller or symmetric case if possible, then finish with a short proof or construction.


\textbf{Common traps.}
\begin{itemize}
\item Brute-forcing cases before identifying the invariant.
\item Assuming symmetry that the rules do not actually provide.
\item Giving a construction without proving that it is minimal or guaranteed.
\end{itemize}


\textbf{Related follow-ups.}
\begin{itemize}
\item Can you generalize the argument to n objects or a different threshold?
\item Is there a second solution using an invariant, parity, or information-theoretic bound?
\end{itemize}


\end{EditorialNote}
\QuestionSpacer

\section{Calculus and Linear Algebra}

\ChapterIntro{Use these questions to practice complete solutions, concise explanations, and likely follow-up questions.}

\subsection{Important Calculus Methods}

\FoundationQuestionHeader{Question 38}{Root-Finding Algorithms}{Important Calculus Methods}

\textbf{Problem.} Explain several numerical root-finding algorithms for solving f(x)=0, including their convergence tradeoffs.


\textbf{Solution.}

Three standard choices are:
\begin{itemize}
\item \textbf{Bisection:} maintain a bracket \([a,b]\) with opposite signs and halve it each step. It is robust and converges linearly.
\item \textbf{Newton:} \(x_{n+1}=x_n-f(x_n)/f'(x_n)\). Near a simple root it converges quadratically, but it needs derivatives and a reasonable starting point.
\item \textbf{Secant:} replace the derivative by a two-point slope. It needs no analytic derivative and has superlinear convergence, but is less robust than bisection.
\end{itemize}
In production, a safeguarded hybrid such as Brent's method is often preferable.


\begin{EditorialNote}


\textbf{Core skill tested.} This problem tests whether foundational calculus or linear algebra can be turned into a compact derivation rather than recalled as an isolated formula.


\textbf{How to present it in interview.} Write the governing identity or transformation first, keep notation clean, and name the numerical-stability or regularity caveat when relevant.


\textbf{Common traps.}
\begin{itemize}
\item Losing a domain or regularity condition.
\item Using a numerically fragile algebraic formula when a factorization is available.
\item Skipping the check that the derived solution satisfies the original equation.
\end{itemize}


\textbf{Related follow-ups.}
\begin{itemize}
\item How would you verify the result numerically?
\item What changes in higher dimensions or when the matrix or function is poorly conditioned?
\end{itemize}


\end{EditorialNote}
\QuestionSpacer

\subsection{Linear Algebra}

\FoundationQuestionHeader{Question 39}{Least Squares via QR}{Linear Algebra}

\textbf{Problem.} If your programming environment has no built-in linear least-squares routine, how would you implement least squares using matrix factorizations?


\textbf{Solution.}

Given a design matrix \(X\) and response vector \(y\), compute a thin QR decomposition
\[
X=QR,
\]
where \(Q^TQ=I\) and \(R\) is upper triangular. Least squares minimizes \(\|X\beta-y\|_2\), which is equivalent to solving
\[
R\hat\beta=Q^Ty.
\]
Use back substitution for \(\hat\beta\). This is numerically preferable to explicitly forming \((X^TX)^{-1}X^Ty\), because the normal equations square the condition number and amplify numerical error.


\begin{EditorialNote}


\textbf{Core skill tested.} This problem tests whether foundational calculus or linear algebra can be turned into a compact derivation rather than recalled as an isolated formula.


\textbf{How to present it in interview.} Write the governing identity or transformation first, keep notation clean, and name the numerical-stability or regularity caveat when relevant.


\textbf{Common traps.}
\begin{itemize}
\item Losing a domain or regularity condition.
\item Using a numerically fragile algebraic formula when a factorization is available.
\item Skipping the check that the derived solution satisfies the original equation.
\end{itemize}


\textbf{Related follow-ups.}
\begin{itemize}
\item How would you verify the result numerically?
\item What changes in higher dimensions or when the matrix or function is poorly conditioned?
\end{itemize}


\end{EditorialNote}
\QuestionSpacer

\section{Probability Theory}

\ChapterIntro{Use these questions to practice complete solutions, concise explanations, and likely follow-up questions.}

\subsection{Basic Probability Definitions and Set Operations}

\FoundationQuestionHeader{Question 40}{Coin Toss Game}{Basic Probability Definitions and Set Operations}

\textbf{Problem.} Gambler A flips n+1 fair coins and gambler B independently flips n fair coins. What is the probability that A obtains strictly more heads than B?


\textbf{Solution.}

The probability is \(\boxed{1/2}\). Couple A's first \(n\) coins with B's \(n\) coins. By symmetry, A is ahead after those \(n\) paired coins as often as B is ahead. If the paired totals tie, A's extra coin breaks the tie in A's favor with probability \(1/2\). A direct binomial calculation gives the same result: for independent \(X\sim\mathrm{Bin}(n+1,1/2)\) and \(Y\sim\mathrm{Bin}(n,1/2)\), \(P(X>Y)=1/2\).


\begin{EditorialNote}


\textbf{Core skill tested.} This problem tests probability modeling: define the sample space or random variables correctly, then use conditioning, symmetry, indicators, or order statistics efficiently.


\textbf{How to present it in interview.} Say what is random and what is conditioned on before calculating. Prefer a one-line symmetry or indicator argument over long enumeration when possible.


\textbf{Common traps.}
\begin{itemize}
\item Mixing ordered and unordered outcomes.
\item Conditioning on an event without renormalizing.
\item Treating dependent events as independent or ignoring a sampling protocol.
\end{itemize}


\textbf{Related follow-ups.}
\begin{itemize}
\item Can you derive the same result with indicators or a complement?
\item How does the answer change under biased or dependent sampling?
\end{itemize}


\end{EditorialNote}
\QuestionSpacer

\subsection{Conditional Probability and Bayes' Formula}

\FoundationQuestionHeader{Question 41}{Double-Headed Coin Posterior}{Conditional Probability and Bayes' Formula}

\textbf{Problem.} Among 1,000 coins, one is double-headed and the other 999 are fair. Choose one coin uniformly at random and observe 10 consecutive heads. What is the posterior probability that you chose the double-headed coin?


\textbf{Solution.}

Let \(U\) be the double-headed coin and \(H^{10}\) the event of ten heads. Bayes' rule gives
\[
P(U\mid H^{10})
=\frac{(1)(1/1000)}{(1)(1/1000)+(2^{-10})(999/1000)}
=\boxed{\frac{1024}{2023}}\approx0.5062.
\]


\begin{EditorialNote}


\textbf{Core skill tested.} This problem tests probability modeling: define the sample space or random variables correctly, then use conditioning, symmetry, indicators, or order statistics efficiently.


\textbf{How to present it in interview.} Say what is random and what is conditioned on before calculating. Prefer a one-line symmetry or indicator argument over long enumeration when possible.


\textbf{Common traps.}
\begin{itemize}
\item Mixing ordered and unordered outcomes.
\item Conditioning on an event without renormalizing.
\item Treating dependent events as independent or ignoring a sampling protocol.
\end{itemize}


\textbf{Related follow-ups.}
\begin{itemize}
\item Can you derive the same result with indicators or a complement?
\item How does the answer change under biased or dependent sampling?
\end{itemize}


\end{EditorialNote}
\QuestionSpacer

\FoundationQuestionHeader{Question 42}{Monty Hall}{Conditional Probability and Bayes' Formula}

\textbf{Problem.} In the standard Monty Hall game, one of three doors hides a car and the other two hide goats. After you choose a door, the host---who knows the prizes---opens a different door showing a goat and offers you the chance to switch to the remaining closed door. Should you switch, and what is your probability of winning if you do?


\textbf{Solution.}

Switch. Your initial door is correct with probability \(1/3\) and wrong with probability \(2/3\). Monty always reveals a goat among the other doors. If your initial choice was wrong, the only remaining unopened door is the car; if it was right, switching loses. Therefore
\[
\boxed{P(\text{win by switching})=\frac23}.
\]


\begin{EditorialNote}


\textbf{Core skill tested.} This problem tests probability modeling: define the sample space or random variables correctly, then use conditioning, symmetry, indicators, or order statistics efficiently.


\textbf{How to present it in interview.} Say what is random and what is conditioned on before calculating. Prefer a one-line symmetry or indicator argument over long enumeration when possible.


\textbf{Common traps.}
\begin{itemize}
\item Mixing ordered and unordered outcomes.
\item Conditioning on an event without renormalizing.
\item Treating dependent events as independent or ignoring a sampling protocol.
\end{itemize}


\textbf{Related follow-ups.}
\begin{itemize}
\item Can you derive the same result with indicators or a complement?
\item How does the answer change under biased or dependent sampling?
\end{itemize}


\end{EditorialNote}
\QuestionSpacer

\subsection{Expected Value, Variance \& Covariance}

\FoundationQuestionHeader{Question 43}{Coupon Collection}{Expected Value, Variance \& Covariance}

\textbf{Problem.} There are N equally likely coupon types, independently drawn. (A) What is the expected number of draws needed to collect all N types? (B) After n draws, what is the expected number of distinct types collected?


\textbf{Solution.}

Suppose there are \(N\) equally likely coupon types.

\textbf{A. Time to complete the set.} After \(i-1\) distinct types have been collected, the probability the next coupon is new is \((N-i+1)/N\), so the expected waiting time for the next new type is \(N/(N-i+1)\). Therefore
\[
\boxed{E[T]=N\sum_{k=1}^{N}\frac1k=NH_N}.
\]

\textbf{B. Distinct types after \(n\) draws.} Let \(I_j\) indicate that type \(j\) has appeared. Then
\[
P(I_j=1)=1-\left(1-\frac1N\right)^n.
\]
By linearity of expectation,
\[
\boxed{E[D_n]=N\left[1-\left(1-\frac1N\right)^n\right]}.
\]


\begin{EditorialNote}


\textbf{Core skill tested.} This problem tests probability modeling: define the sample space or random variables correctly, then use conditioning, symmetry, indicators, or order statistics efficiently.


\textbf{How to present it in interview.} Say what is random and what is conditioned on before calculating. Prefer a one-line symmetry or indicator argument over long enumeration when possible.


\textbf{Common traps.}
\begin{itemize}
\item Mixing ordered and unordered outcomes.
\item Conditioning on an event without renormalizing.
\item Treating dependent events as independent or ignoring a sampling protocol.
\end{itemize}


\textbf{Related follow-ups.}
\begin{itemize}
\item Can you derive the same result with indicators or a complement?
\item How does the answer change under biased or dependent sampling?
\end{itemize}


\end{EditorialNote}
\QuestionSpacer

\section{Stochastic Process and Stochastic Calculus}

\ChapterIntro{Use these questions to practice complete solutions, concise explanations, and likely follow-up questions.}

\subsection{Markov Chain}

\FoundationQuestionHeader{Question 44}{Gambler's Ruin -- Unequal Win Probability}{Markov Chain}

\textbf{Problem.} Player \(M\) starts with \$1 and player \(N\) with \$2. In each game, the winner receives \$1 from the loser. \(M\) wins each game independently with probability \(2/3\). Play continues until one player is bankrupt. What is the probability that \(M\) wins all \$3?


\textbf{Solution.}

Let \(p_i\) be the probability that M eventually owns all \$3 when M currently owns \(i\) dollars. Then \(p_0=0\), \(p_3=1\), and because M wins each game with probability \(2/3\),
\[
p_1=\frac23p_2,\qquad
p_2=\frac23+\frac13p_1.
\]
Solving,
\[
p_1=\frac23\left(\frac23+\frac13p_1\right)
=\frac49+\frac29p_1,
\]
so
\[
\boxed{p_1=\frac47}.
\]
This is the gambler's-ruin absorption probability starting from M's initial wealth of \$1.


\begin{EditorialNote}


\textbf{Core skill tested.} This problem tests stochastic-process state design, stopping-time reasoning, or the use of martingales/Markov structure to avoid unnecessary path enumeration.


\textbf{How to present it in interview.} Define the state or stopping time, write the recursion/PDE/martingale, then state boundary conditions before solving.


\textbf{Common traps.}
\begin{itemize}
\item Choosing a state that forgets information needed for the future.
\item Applying optional stopping without checking the setup.
\item Confusing an unconditional mean-zero process with a martingale.
\end{itemize}


\textbf{Related follow-ups.}
\begin{itemize}
\item Can you solve it both with a Markov recursion and a martingale argument?
\item What changes with drift, bias, or different absorbing boundaries?
\end{itemize}


\end{EditorialNote}
\QuestionSpacer

\subsection{Dynamic Programming}

\FoundationQuestionHeader{Question 45}{Three-Roll Optimal Stopping}{Dynamic Programming}

\textbf{Problem.} You may roll a fair six-sided die at most three times. After the first or second roll, you may either accept that roll's face value in dollars and stop, or discard it and continue. The third roll, if reached, must be accepted. What is the optimal strategy and fair value of the game?


\textbf{Solution.}

Work backward.

On the third roll there is no choice, so its value is
\[
V_3=E[X]=3.5.
\]
After the second roll, keep \(4,5,6\) and reroll \(1,2,3\). Therefore
\[
V_2=\frac{3(3.5)+4+5+6}{6}=\frac{25.5}{6}=4.25.
\]
After the first roll, keep \(5\) or \(6\), and reroll \(1,2,3,4\), because continuation is worth 4.25. Hence
\[
V_1=\frac{4(4.25)+5+6}{6}=\boxed{\frac{14}{3}}\approx4.667.
\]
Optimal policy: after roll 1 stop on 5--6; after roll 2 stop on 4--6; on roll 3 accept whatever appears.


\begin{EditorialNote}


\textbf{Core skill tested.} This problem tests stochastic-process state design, stopping-time reasoning, or the use of martingales/Markov structure to avoid unnecessary path enumeration.


\textbf{How to present it in interview.} Define the state or stopping time, write the recursion/PDE/martingale, then state boundary conditions before solving.


\textbf{Common traps.}
\begin{itemize}
\item Choosing a state that forgets information needed for the future.
\item Applying optional stopping without checking the setup.
\item Confusing an unconditional mean-zero process with a martingale.
\end{itemize}


\textbf{Related follow-ups.}
\begin{itemize}
\item Can you solve it both with a Markov recursion and a martingale argument?
\item What changes with drift, bias, or different absorbing boundaries?
\end{itemize}


\end{EditorialNote}
\QuestionSpacer

\section{Finance}

\ChapterIntro{Use these questions to practice complete solutions, concise explanations, and likely follow-up questions.}

\subsection{Option Pricing}

\FoundationQuestionHeader{Question 46}{Put--Call Parity}{Option Pricing}

\textbf{Problem.} State and prove put-call parity for European options on a non-dividend-paying stock.


\textbf{Solution.}

For a non-dividend-paying stock, compare two portfolios with maturity \(T\):
\begin{itemize}
\item one European call plus cash worth \(K\) at maturity;
\item one European put plus one share of stock.
\end{itemize}
Their terminal payoffs are identical:
\[
(S_T-K)^++K=(K-S_T)^++S_T=\max(S_T,K).
\]
No-arbitrage therefore implies
\[
\boxed{C+Ke^{-rT}=P+S_0},
\]
or
\[
\boxed{C-P=S_0-Ke^{-rT}}.
\]
With a continuous dividend yield \(q\), replace \(S_0\) by \(S_0e^{-qT}\).


\begin{EditorialNote}


\textbf{Core skill tested.} This problem tests no-arbitrage reasoning, replication, risk-neutral valuation, or disciplined interpretation of option/rate sensitivities.


\textbf{How to present it in interview.} Start from payoff replication or the pricing measure, write the minimal formula, and immediately state the assumptions under which the sign or equality holds.


\textbf{Related original-bank record.} \texttt{a related question}.


\textbf{Common traps.}
\begin{itemize}
\item Quoting a Greek or formula without its carry/dividend assumptions.
\item Treating a local hedge as a risk-free arbitrage.
\item Ignoring convexity, tail risk, units, or discounting conventions.
\end{itemize}


\textbf{Related follow-ups.}
\begin{itemize}
\item Can you derive the result from both replication and risk-neutral valuation?
\item Which assumption is most likely to fail on a real desk?
\end{itemize}


\end{EditorialNote}
\QuestionSpacer

\subsection{The Greeks}

\FoundationQuestionHeader{Question 47}{Dynamic Delta Hedging}{The Greeks}

\textbf{Problem.} You are long a European call and want to delta-hedge it with the underlying stock. What initial stock position should you take, and how should you rebalance after a sudden increase in the stock price?


\textbf{Solution.}

A long call has positive delta. To make the position locally delta-neutral, short
\[
\boxed{\Delta_C\text{ shares per call}}
\]
of stock (with the appropriate contract multiplier).

A vanilla long call has positive gamma, so if GM rises, the call delta increases and the existing short-stock hedge is no longer large enough; short additional shares. If GM falls, delta decreases; buy back some of the short stock.

This is dynamic delta hedging. In continuous Black--Scholes it removes infinitesimal spot risk, but in practice discrete rebalancing leaves gamma, jump, transaction-cost, and model risk.


\begin{EditorialNote}


\textbf{Core skill tested.} This problem tests no-arbitrage reasoning, replication, risk-neutral valuation, or disciplined interpretation of option/rate sensitivities.


\textbf{How to present it in interview.} Start from payoff replication or the pricing measure, write the minimal formula, and immediately state the assumptions under which the sign or equality holds.


\textbf{Common traps.}
\begin{itemize}
\item Quoting a Greek or formula without its carry/dividend assumptions.
\item Treating a local hedge as a risk-free arbitrage.
\item Ignoring convexity, tail risk, units, or discounting conventions.
\end{itemize}


\textbf{Related follow-ups.}
\begin{itemize}
\item Can you derive the result from both replication and risk-neutral valuation?
\item Which assumption is most likely to fail on a real desk?
\end{itemize}


\end{EditorialNote}
\QuestionSpacer

\subsection{Other Finance Questions}

\FoundationQuestionHeader{Question 48}{Value at Risk}{Other Finance Questions}

\textbf{Problem.} Define Value at Risk (VaR). What are important shortcomings of VaR, especially for nonlinear derivative portfolios?


\textbf{Solution.}

Value at Risk at confidence level \(\alpha\) and horizon \(h\) is a loss quantile: the threshold \(v\) such that losses exceed \(v\) with probability about \(1-\alpha\). In notation for a loss variable \(L\),
\[
\operatorname{VaR}_\alpha(L)=\inf\{\ell:P(L\le\ell)\ge\alpha\}.
\]

Major limitations are especially important for derivatives:
\begin{itemize}
\item VaR says little about how large losses are beyond the quantile;
\item nonlinear payoffs and volatility/skew make normal approximations unreliable;
\item estimates are sensitive to model, horizon, and historical window;
\item VaR is not subadditive in full generality, so it can reward apparent diversification incorrectly.
\end{itemize}
Expected shortfall is often preferred for tail severity because it averages losses beyond VaR and is coherent under standard definitions.


\begin{EditorialNote}


\textbf{Core skill tested.} This problem tests no-arbitrage reasoning, replication, risk-neutral valuation, or disciplined interpretation of option/rate sensitivities.


\textbf{How to present it in interview.} Start from payoff replication or the pricing measure, write the minimal formula, and immediately state the assumptions under which the sign or equality holds.


\textbf{Common traps.}
\begin{itemize}
\item Quoting a Greek or formula without its carry/dividend assumptions.
\item Treating a local hedge as a risk-free arbitrage.
\item Ignoring convexity, tail risk, units, or discounting conventions.
\end{itemize}


\textbf{Related follow-ups.}
\begin{itemize}
\item Can you derive the result from both replication and risk-neutral valuation?
\item Which assumption is most likely to fail on a real desk?
\end{itemize}


\end{EditorialNote}
\QuestionSpacer

\section{Algorithms and Numerical Methods}

\ChapterIntro{Use these questions to practice complete solutions, concise explanations, and likely follow-up questions.}

\subsection{Algorithms}

\FoundationQuestionHeader{Question 49}{Fisher--Yates Shuffle}{Algorithms}

\textbf{Problem.} Given access to a uniform random-number generator, shuffle a 52-card deck so that every permutation is equally likely. Give the algorithm and justify its correctness.


\textbf{Solution.}

Use the Fisher--Yates shuffle. For \(i=51,50,\ldots,1\), draw
\[
j\sim\operatorname{Unif}\{0,1,\ldots,i\}
\]
and swap positions \(i\) and \(j\).

At step \(i\), each of the \(i+1\) remaining cards is equally likely to be placed in position \(i\). Multiplying the conditional probabilities shows every one of the \(52!\) permutations has probability
\[
\boxed{1/52!}.
\]
The runtime is \(O(52)\), or \(O(n)\) for \(n\) items.


\begin{EditorialNote}


\textbf{Core skill tested.} This problem tests algorithm design: identify the invariant/data structure, prove correctness, and meet the stated time/space target.


\textbf{How to present it in interview.} Give the invariant and complexity before code. Then show the update/search logic and mention edge cases such as empty input, overflow, or unknown bounds.


\textbf{Common traps.}
\begin{itemize}
\item Giving code without a correctness invariant.
\item Hiding an O(n) operation inside an apparently constant-time loop.
\item Ignoring overflow, boundary conditions, or worst-case behavior.
\end{itemize}


\textbf{Related follow-ups.}
\begin{itemize}
\item Can the solution be made streaming or in-place?
\item What are the worst-case time and memory bounds, and how would you test them?
\end{itemize}


\end{EditorialNote}
\QuestionSpacer

\subsection{Numerical Methods}

\FoundationQuestionHeader{Question 50}{Monte Carlo Variance Reduction}{Numerical Methods}

\textbf{Problem.} Explain several variance-reduction techniques that can improve the efficiency of Monte Carlo pricing.


\textbf{Solution.}

Common variance-reduction techniques include:
\begin{itemize}
\item \textbf{Antithetic variates:} pair a path driven by \(Z\) with one driven by \(-Z\); averaging helps when the payoffs are negatively correlated.
\item \textbf{Control variates:} use a correlated quantity \(Y\) with known expectation and estimate \(X-\beta(Y-E[Y])\), choosing \(\beta\) near the regression-optimal coefficient.
\item \textbf{Importance sampling:} simulate from a distribution that visits the payoff-relevant region more often and reweight by the likelihood ratio; especially useful for rare events/deep OTM payoffs.
\item \textbf{Stratified or quasi-Monte Carlo sampling:} cover the input space more evenly than independent pseudorandom draws.
\item \textbf{Moment matching:} adjust a simulated sample so low-order moments better match the target distribution, though this should be used carefully to avoid hidden bias.
\end{itemize}
The best method depends on the payoff and on which source of variance dominates.


\begin{EditorialNote}


\textbf{Core skill tested.} This problem tests algorithm design: identify the invariant/data structure, prove correctness, and meet the stated time/space target.


\textbf{How to present it in interview.} Give the invariant and complexity before code. Then show the update/search logic and mention edge cases such as empty input, overflow, or unknown bounds.


\textbf{Common traps.}
\begin{itemize}
\item Giving code without a correctness invariant.
\item Hiding an O(n) operation inside an apparently constant-time loop.
\item Ignoring overflow, boundary conditions, or worst-case behavior.
\end{itemize}


\textbf{Related follow-ups.}
\begin{itemize}
\item Can the solution be made streaming or in-place?
\item What are the worst-case time and memory bounds, and how would you test them?
\end{itemize}


\end{EditorialNote}
\QuestionSpacer
\end{document}
